In this section, we assume any vector space $V$ is finite-dimensional over $\mathbb{R}.$

\subsection{Bilinear and Quadratic Forms}

\begin{definition}[Bilinear Form]
We say a function $\tau: V\times V\to\mathbb{R}$ is bilinear if it is linear in each component fixing the other one.
i.e.(Exercise)
\end{definition}

\begin{remark}
    Real inner product space is a special case of bilinear form with additonal assumptions (exercise). 
\end{remark}

\begin{definition}
    Let $\beta=\{\mathbf{v}_1,\mathbf{v}_2,\cdots,\mathbf{v}_n\}$ be a basis of $V$ and $\tau:V\times V\to\mathbb{R}$ be a bilinear form. The matrix representing a bilinear form w.r.t $\beta$ is given by a square matrix $\psi_\beta(\tau)=A=(A_{ij})$ by setting $A_{ij}=\tau(v_i,v_j)$.
\end{definition}

\begin{theorem}
     Let $\beta$ be a basis of $V$ and $\tau:V\times V\to\mathbb{R}$ be a bilinear form, for any $\mathbf{v},\mathbf{w}\in V$, one has:
     \[\tau(\mathbf{v},\mathbf{w})=[\mathbf{v}]_\beta^\top\psi_\beta(\tau)[\mathbf{w}]_\beta.\]
     \label{bilinear matrix vector product}
\end{theorem}

\begin{pfbox}
    It suffices to verify the above equality for any two basis vectors.
\end{pfbox}

\begin{proposition}
    Let $\tau:V\times V\to\mathbb{R}$ be a non-zero bilinear form, TFAE:
    \begin{enumerate}
        \item $\rank(\psi_\beta(\tau))=1$ for any basis $\beta$ of $V$.
        \item There exist two linear functionals $f,g\in V^*$ such that $\tau(\mathbf{v},\mathbf{w})=f(\mathbf{v})g(\mathbf{w}).$
    \end{enumerate}
    \label{linear map product and bilinear}
\end{proposition}

\begin{corollary}
    Product of two linear functionals is a bilinear form of at most rank one.
\end{corollary}

\begin{remark}
    The converse is not true. Consider the real dot product and use~\eqref{linear map product and bilinear} to disprove it. 
\end{remark}


\begin{theorem}[Congruence]
    Let $\beta_1, \beta_2$ be two bases of $V$ and $\tau: V\times V\to \mathbb{R}$ be a bilinear form.
    Set $A=\psi_{\beta_1}(\tau)$ and $B=\psi_{\beta_2}(\tau)$, then there exists an invertible matrix $P$ such that $A=P^\top BP.$ 
\end{theorem}

\begin{pfbox}
    (Sketch) Recall $P$ is a change of coordinate matrix from $\beta_1$ to $\beta_2$, then for any $\mathbf{v}\in V$, $P[\mathbf{v}]_{\beta_1} = [\mathbf{v}]_{\beta_2}.$
    Use theorem~\eqref{bilinear matrix vector product} to deduce.
\end{pfbox}


\begin{definition}
Two matrices $A,B\in M_n(\mathbb{R})$ are congruent if there exists an invertible matrix $P$ such that
\[
A=P^\top BP.
\]
\label{congruence}
\end{definition}

\begin{remark}\ \vspace{-1em}
    \begin{enumerate}
        \item For linear operators, we say two matrices $A$ and $B$ are similar if...
        \item For the congruence of bilinear forms, $P^\top$ is not necessarily $P^{-1}.$ 
    \end{enumerate}
\end{remark}

\begin{definition}
    A bilinear form is called
    \begin{enumerate}
        \item symmetric if $\tau(\mathbf{v}_1,\mathbf{v}_2)=\tau(\mathbf{v}_2,\mathbf{v}_1)$ for any $\mathbf{v}_1,\mathbf{v}_2\in V$.
        \item anti-symmetric if $\tau(\mathbf{v}_1,\mathbf{v}_2)=-\tau(\mathbf{v}_2,\mathbf{v}_1)$ for any $\mathbf{v}_1,\mathbf{v}_2\in V$.
    \end{enumerate}
\end{definition}

\begin{theorem}
    Fix a basis $\beta$ of $V$
    \begin{enumerate}
        \item $\tau$ is symmetric iff $\psi_\beta(\tau)$ is a symmetric matrix.
        \item $\tau$ is anti-symmetric iff $\psi_\beta(\tau)$ is an anti-symmetric matrix.
    \end{enumerate}
\end{theorem}

\begin{pfbox}
    Omit.
\end{pfbox}
\

\begin{theorem}
    Any bilinear form $\tau$ can be uniquely written as a sum $\tau_1+\tau_2,$ where $\tau_1$ is symmetric and $\tau_2$ is anti-symmetric.
\end{theorem}

\begin{pfbox}
    First observe that any bilinear form 
    \[\tau(\mathbf{v}, \mathbf{w})=\underbrace{\frac{\tau(\mathbf{v},\mathbf{w})+\tau(\mathbf{w},\mathbf{v})}{2}}_{\text{Symmetric}} + \underbrace{\frac{\tau(\mathbf{v},\mathbf{w})-\tau(\mathbf{w},\mathbf{v})}{2}}_{\text{Anti-symmetric}}. \]
    For the uniqueness, suppose $\tau=\tau_1+\tau_2=\tau'_1+\tau'_2,$ then $\tau_1-\tau'_1 = \tau'_2-\tau_2 \equiv \delta.$ Note that $\delta$ is both symmetric and anti-symmetric, this forces $\delta\equiv 0.$
\end{pfbox}



\begin{definition}[Quadratic form]
A quadratic form $q:V\to\mathbb{R}$ is a function of the form
\[
q(\mathbf{v})=\tau(\mathbf{v},\mathbf{v})
\]
for some bilinear form $\tau:V\times V\to\mathbb{R}$.
\end{definition}

\begin{theorem}
There is a one-to-one correspondence between:
\begin{enumerate}
    \item Quadratic forms $q:V\to\mathbb{R}$.
    \item Symmetric bilinear forms $\tau:V\times V\to\mathbb{R}$.
    \item Symmetric matrices $A\in \operatorname{Sym}_n(\mathbb{R})$.
\end{enumerate}
\end{theorem}

\begin{pfbox}
We know the bijection between symmetric bilinear forms and symmetric matrices via $\psi_\beta(\tau)$.

To show the bijection between (1) and (2), given a symmetric bilinear form $\tau$,
$q(\mathbf{v})=\tau(\mathbf{v},\mathbf{v})$ is clearly a quadratic form.

Conversely, given any quadratic form $q$, we have $q(\mathbf{v})=\tau(\mathbf{v},\mathbf{v})$ for some bilinear form $\tau$.
Note that $\tau(\mathbf{v},\mathbf{v})=\tau_1(\mathbf{v},\mathbf{v})+\underbrace{\tau_2(\mathbf{v},\mathbf{v})}_{0}$,
where $\tau_1$ is symmetric and $\tau_2$ is anti-symmetric.
Via the Polarization identity, we can recover $\tau_1$ from $q$:
\[
\tau_1(\mathbf{v},\mathbf{w})
=
\frac{q(\mathbf{v}+\mathbf{w})-q(\mathbf{v})-q(\mathbf{w})}{2}.
\]
This identity establishes a unique symmetric bilinear form for every quadratic form.
\end{pfbox}

\begin{remark}
If $V=\mathbb{R}^n$, the quadratic form coincides with a homogeneous polynomial of degree $2$:
\[
q(\mathbf{x})
=
\mathbf{x}^\top A\mathbf{x}
=
\sum_{i=1}^n \alpha_{ii}x_i^2
+
\sum_{i<j}2\alpha_{ij}x_ix_j,
\]
where $A\in \operatorname{Sym}_n(\mathbb{R})$.
\end{remark}


\begin{theorem}[Diagonalization of symmetric bilinear forms]
Let $(V,\langle -,-\rangle)$ be a finite-dimensional real inner product space. Then every symmetric bilinear form $\tau$ is diagonalizable.
\label{Diagonalization of symmetric bilinear forms}
\end{theorem}

\begin{pfbox}
Fix $\mathbf{v}\in V$. Consider the linear functional
$\tau: V\to\mathbb{R}$ defined by $\mathbf{w}\mapsto\tau(\mathbf{v},\mathbf{w})\,\,\forall\,\mathbf{w}\in V$.
By the Riesz Representation Theorem, there exists a unique vector $T(\mathbf{v})\in V$ such that
$\tau(\mathbf{v},\mathbf{w})=\innerproduct{T(\mathbf{v})}{\mathbf{w}}$
for all $\mathbf{w}\in V$.

This defines an operator $T:V\to V$.
Because $\tau$ is bilinear, $T$ is linear. Since $\tau$ is symmetric,
\[
\innerproduct{T\mathbf{v}}{\mathbf{w}}
=
\tau(\mathbf{v},\mathbf{w})
=
\tau(\mathbf{w},\mathbf{v})
=
\innerproduct{T\mathbf{w}}{\mathbf{v}}
=
\innerproduct{\mathbf{v}}{T\mathbf{w}},
\]
so $T$ is self-adjoint (Hermitian).

By the Real Spectral Theorem, there exists an ONB $\beta=\{\mathbf{v}_1,\dots,\mathbf{v}_n\}$ of eigenvectors of $T$ such that
\[
\tau(\mathbf{v}_i,\mathbf{v}_j)
=
\innerproduct{T\mathbf{v}_i}{\mathbf{v}_j}
=
\innerproduct{\lambda_i\mathbf{v}_i}{\mathbf{v}_j}
=
\lambda_i\delta_{ij}.
\]
Hence $\psi_\beta(\tau)$ is a diagonal matrix.
\end{pfbox}

\begin{remark}
    Spectral theorem gives an ONB that helps to diagonalize $\tau$, but it does not rule out the possibility that there could be any other choice of basis.
    Indeed, the diagonalization of a bilinear form is not unique.
\end{remark}


\begin{theorem}[Diagonalization of Quadratic Forms]
Let $(V,\langle -,-\rangle)$ be an inner product space and $q:V\to\mathbb{R}$ be a quadratic form. Then there exists a basis $\{\mathbf{v}_1,\dots,\mathbf{v}_n\}$ such that
\[
q(\mathbf{v})
=
\sum_{i=1}^n \alpha_i x_i^2,
\]
where
\[
\mathbf{v}=\sum_{i=1}^n x_i\mathbf{v}_i
\qquad\text{and}\qquad
\alpha_i\in\mathbb{R}.
\]
\end{theorem}

\begin{pfbox}
By the diagonalization of symmetric bilinear forms, there is a basis\\
$\beta=\{\mathbf{v}_1,\dots,\mathbf{v}_n\}$ such that
$\tau(\mathbf{v}_i,\mathbf{v}_j)=\alpha_i\delta_{ij}$.
Then
\begin{align*}
q(\mathbf{v})
&=
\tau\left(\sum x_i\mathbf{v}_i,\sum x_j\mathbf{v}_j\right) \\
&=
\sum_{i=1}^n x_i^2\tau(\mathbf{v}_i,\mathbf{v}_i) = \sum_{i=1}^n \alpha_i x_i^2.
\end{align*}
\end{pfbox}

\begin{remark}
    Diagonaliability of quadratic forms confirms that completing the square of any quadratic form is not unique.
\end{remark}

\begin{theorem}[Diagonalization of Symmetric Matrix]
Let $A\in \operatorname{Sym}_n(\mathbb{R})$. Then $A$ is congruent to a diagonal matrix.
\end{theorem}
\begin{pfbox}
    Direct result from the definition~\eqref{congruence} and theorem~\eqref{Diagonalization of symmetric bilinear forms}
\end{pfbox}

\begin{definition}
Let $A\in \operatorname{Sym}_n(\mathbb{R})$, and $D$ be a diagonal matrix congruent to $A$.
Let $t$ be the number of positive diagonal entries of $D$, and $m$ be the number of negative diagonal entries of $D$.
Then 
\begin{enumerate}
    \item $t$ is called the \textbf{index} of $A$,
    \item $t-m$ is called the \textbf{signature} of $A$.
\end{enumerate}
\end{definition}

\begin{theorem}[Sylvester's Law of Inertia]
Let $A\in \operatorname{Sym}_n(\mathbb{R})$. Suppose $A$ is congruent to two diagonal matrices $D_1$ and $D_2$. Then $D_1$ and $D_2$ share the same index and signature.
\end{theorem}

\begin{pfbox}
Let $\beta_1=\{\mathbf{v}_1,\dots,\mathbf{v}_n\},
\,
\beta_2=\{\mathbf{v}_1',\dots,\mathbf{v}_n'\}$
such that
\[
\psi_{\beta_1}(\tau)
=
D_1
=
\operatorname{diag}(\alpha_1,\dots,\alpha_n)
\]
and
\[
\psi_{\beta_2}(\tau)
=
D_2
=
\operatorname{diag}(\alpha_1',\dots,\alpha_n').
\]
Define $t_1,m_1$ as the number of positive/negative entries in $D_1$, and $t_2,m_2$ for $D_2$.

Note that
\[
\rank(D_1)
=
 t_1+m_1
=
\rank(A)
=
 t_2+m_2
=
\rank(D_2).
\]
It suffices to show $t_1=t_2$.

Suppose, WLOG, $t_1>t_2$. Assume the first $t_1$ and $t_2$ elements are positive. Consider
$V_1=\spn\{\mathbf{v}_1,\dots,\mathbf{v}_{t_1}\}$ and
$V_2'=\spn\{\mathbf{v}_{t_2+1}',\dots,\mathbf{v}_n'\}$.

\textbf{Claim.} $V_1\cap V_2'=\{\mathbf{0}\}$.
Suppose not. Take $\mathbf{v}\in V_1\cap V_2'$ such that $\mathbf{v}\ne \mathbf{0}$. Since $\mathbf{v}\in V_1$,
\[
q(\mathbf{v})
=
\sum_{i=1}^{t_1}\alpha_i x_i^2
>
0,
\]
because $\alpha_i>0$ and $\mathbf{v}\ne \mathbf{0}$. Since $\mathbf{v}\in V_2'$,
\[
q(\mathbf{v})
=
\sum_{i=t_2+1}^n \alpha_i' y_i^2
\le
0,
\]
because $\alpha_i'\le 0$. This is absurd.

On the other hand, the dimension formula gives
\begin{align*}
\dim(V_1\cap V_2')
&=
\dim(V_1)+\dim(V_2')-\dim(V_1+V_2') \\
&=
 t_1+(n-t_2)-\dim(V_1+V_2') \\
&\ge
 t_1+n-t_2-n \\
&>
0.
\end{align*}
This contradicts the claim. So we must have $t_1=t_2$.
\end{pfbox}


\begin{theorem}
Let $A\in \operatorname{Sym}_n(\mathbb{R})$ and $P$ be any invertible matrix. The number of positive eigenvalues of $A$ and $P^\top AP$ are equal.
\end{theorem}

\begin{pfbox}
By the spectral theorem, there exists an orthogonal matrix $Q$ such that
\[
Q^\top AQ
=
D_1
=
\operatorname{diag}(\lambda_1,\dots,\lambda_n).
\]
Thus $A$ is congruent to $D_1$. The number of positive eigenvalues is the index $t$.
Let $C=P^\top AP$, then $C^\top=C$ is symmetric. By the spectral theorem, there exists an orthogonal matrix $\widetilde{Q}$ such that
\[
\widetilde{Q}^\top C\widetilde{Q}
=
D_2
=
\operatorname{diag}(\mu_1,\dots,\mu_n).
\]
So $C$ is congruent to $D_2$. But $C$ is congruent to $A$, which means $A$ is congruent to $D_2$.
By Sylvester's Law of Inertia, $D_1$ and $D_2$ must have the same index. Thus, the number of positive eigenvalues is preserved.
\end{pfbox}

\begin{corollary}
Let $A\in \operatorname{Sym}_n(\mathbb{R})$ and $P$ be any invertible matrix. The number of negative eigenvalues of $A$ and $P^\top AP$ are equal.
\end{corollary}

\begin{corollary}\ \vspace{-2em}
\begin{enumerate}
	\item $A\in \operatorname{Sym}_n(\mathbb{R})$ is PD iff the index of $A$ is $n$.
	\item $A\in \operatorname{Sym}_n(\mathbb{R})$ is ND iff the signature of $A$ is $-n$.
\end{enumerate}
\end{corollary}

\begin{theorem}
Let $A,B\in \operatorname{Sym}_n(\mathbb{R})$. TFAE:
\begin{enumerate}
    \item[(1)] $A$ and $B$ are congruent.
    \item[(2)] $A$ and $B$ share the same index and signature.
\end{enumerate}
\end{theorem}

\begin{pfbox}
\textbf{(1) $\implies$ (2).}
Suppose $A=P^\top BP$.
By the previous theorem and corollary, $A$ and $B$ have the same number of positive and negative eigenvalues, so they share the same index and signature.

\medskip

\textbf{(2) $\implies$ (1).}
Suppose
\[
A\sim D_1=\operatorname{diag}(\lambda_1,\dots,\lambda_n)
\]
and
\[
B\sim D_2=\operatorname{diag}(\mu_1,\dots,\mu_n).
\]
Since they have the same index $t$ and negative eigenvalues $m$, we can assume the first $t$ entries are positive, the next $m$ are negative, and the rest are zero for both $D_1$ and $D_2$.

Consider the diagonal matrix
\[
Q
=
\operatorname{diag}
\left(
\sqrt{\frac{\mu_1}{\lambda_1}},
\dots,
\sqrt{\frac{\mu_t}{\lambda_t}},
\sqrt{\frac{\mu_{t+1}}{\lambda_{t+1}}},
\dots,
1,
\dots
\right).
\]
It is clear that $Q$ is invertible and
\[
D_2=Q^\top D_1Q.
\]
Write
\[
D_1=P^\top AP,
\qquad
D_2=\widetilde{P}^{\top}B\widetilde{P}.
\]
Then
\[
\widetilde{P}^{\top}B\widetilde{P}
=
Q^\top(P^\top AP)Q.
\]
Solving for $B$,
\begin{align*}
B
&=
(\widetilde{P}^{\top})^{-1}Q^\top P^\top APQ\widetilde{P}^{-1} \\
&=
(PQ\widetilde{P}^{-1})^\top A(PQ\widetilde{P}^{-1}).
\end{align*}
Thus, $A$ and $B$ are congruent.
\end{pfbox}


\begin{corollary}
If $A$ is PD/ND/PSD/NSD, then any matrix congruent to $A$ is also PD/ND/PSD/NSD.
\end{corollary}

\begin{pfbox}
For PD/ND case, we can prove them with an alternative method.\\
Let $B$ be congruent to $A$, so there is an invertible $P$ such that $B=P^\top AP$.\\
Let $q_A(\mathbf{v})=\mathbf{v}^\top A\mathbf{v}$, for $B$, $q_B(\mathbf{v})=\mathbf{v}^\top(P^\top AP)\mathbf{v}=(P\mathbf{v})^\top A(P\mathbf{v})$.

Let $\mathbf{w}=P\mathbf{v}$, since $A$ is PD,
$q_A(\mathbf{w})>0 $ for all $\mathbf{w}\ne \mathbf{0}$. 
Since $P$ is invertible, $q_B(\mathbf{v})=q_A(\mathbf{w})>0\,\,\forall\,\mathbf{v}\ne\mathbf{0}$, meaning $B$ is also PD.
\end{pfbox}

\subsection{Definiteness}

\begin{definition}[Definiteness]
	Let $A$ be a symmetric matrix. Then $A$ is called:
	\begin{enumerate}
		\item \textbf{positive-definite (PD)} if $\lambda > 0$ for all eigenvalues $\lambda$ of $A$.
		\item \textbf{positive-semidefinite (PSD)} if $\lambda \ge 0$ for all eigenvalues $\lambda$ of $A$.
		\item \textbf{negative-definite (ND)} if $\lambda < 0$ for all eigenvalues $\lambda$ of $A$.
		\item \textbf{negative-semidefinite (NSD)} if $\lambda \le 0$ for all eigenvalues $\lambda$ of $A$.
		\item \textbf{indefinite} if it is neither PSD nor NSD (i.e., there exists at least one strictly positive eigenvalue and one strictly negative eigenvalue).
	\end{enumerate}
\end{definition}

\begin{remark}
	In particular, PD and ND matrices are special kinds of PSD and NSD matrices respectively. This definition can be extended to Hermitian matrices because all the eigenvalues of a Hermitian matrix are real.
\end{remark}


\begin{theorem}
	Let $A$ be an $n \times n$ symmetric matrix and $q_A(\mathbf{v}) = \mathbf{v}^t A \mathbf{v}$ be its associated quadratic form. Note that clearly $q_A(\mathbf{0}) = 0$.
	\begin{enumerate}
		\item $A$ is PD if and only if $q_A(\mathbf{v}) > 0$ for any $\mathbf{v} \in \mathbb{R}^n \setminus \{\mathbf{0}\}$.
		\item $A$ is PSD if and only if $q_A(\mathbf{v}) \ge 0$ for any $\mathbf{v} \in \mathbb{R}^n \setminus \{\mathbf{0}\}$.
		\item $A$ is ND if and only if $q_A(\mathbf{v}) < 0$ for any $\mathbf{v} \in \mathbb{R}^n \setminus \{\mathbf{0}\}$.
		\item $A$ is NSD if and only if $q_A(\mathbf{v}) \le 0$ for any $\mathbf{v} \in \mathbb{R}^n \setminus \{\mathbf{0}\}$.
		\item $A$ is indefinite if and only if there exist $\mathbf{v}_1, \mathbf{v}_2 \in \mathbb{R}^n$ such that $q_A(\mathbf{v}_1) \le 0$ and $q_A(\mathbf{v}_2) > 0$.
	\end{enumerate}
\end{theorem}

\begin{pfbox}
	As $A$ is symmetric, by the Real Spectral Theorem, there exists an orthonormal basis (ONB) $\{\mathbf{u}_1, \mathbf{u}_2, \dots, \mathbf{u}_n\}$ of $\mathbb{R}^n$ consisting of eigenvectors of $A$, such that the quadratic form can be written as:
	$$
		q_A(\mathbf{v}) = \lambda_1 \innerproduct{\mathbf{v}}{\mathbf{u}_1}^2 + \lambda_2 \innerproduct{\mathbf{v}}{\mathbf{u}_2}^2 + \dots + \lambda_n \innerproduct{\mathbf{v}}{\mathbf{u}_n}^2
	$$
	where $\lambda_i$ is the eigenvalue of $A$ corresponding to $\mathbf{u}_i$. The other parts are similar, so we prove only the first statement (PD case).

	$(\Rightarrow)$ Suppose $A$ is PD. By definition, all eigenvalues satisfy $\lambda_1, \lambda_2, \dots, \lambda_n > 0$. For any $\mathbf{v} \in \mathbb{R}^n \setminus \{\mathbf{0}\}$, since $\{\mathbf{u}_i\}$ forms a basis, we must have $\innerproduct{\mathbf{v}}{\mathbf{u}_i} \ne 0$ for some $i$. Since $\lambda_i > 0$ and $\innerproduct{\mathbf{v}}{\mathbf{u}_j}^2 \ge 0$ for all $j$, it follows that
	$$
		q_A(\mathbf{v}) \ge \lambda_i \innerproduct{\mathbf{v}}{\mathbf{u}_i}^2 > 0.
	$$

	$(\Leftarrow)$ Suppose $q_A(\mathbf{v}) > 0$ for any $\mathbf{v} \in \mathbb{R}^n \setminus \{\mathbf{0}\}$. In particular, choosing $\mathbf{v} = \mathbf{u}_i$ for each $i \in \{1, \dots, n\}$, we have $\innerproduct{\mathbf{u}_i}{\mathbf{u}_j} = \delta_{ij}$. Thus,
	\begin{align*}
		q_A(\mathbf{u}_i)
		 & = \lambda_i \innerproduct{\mathbf{u}_i}{\mathbf{u}_i}^2 \\
		 & = \lambda_i>0.
	\end{align*}
	Since all $\lambda_i > 0$, $A$ is PD by definition.
\end{pfbox}






\begin{definition}[Leading Principal Minors]
	Let $A$ be an $n \times n$ matrix. The leading principal minors (LPM) are the determinants of the upper-left submatrices of $A$:
	\begin{itemize}
		\item $LPM_1$: the determinant of the upper-left $1 \times 1$ corner of $A$,
		\item $LPM_2$: the determinant of the upper-left $2 \times 2$ corner of $A$,
		\item \dots
		\item $LPM_n$: the determinant of $A$ itself (i.e., $\det(A)$).
	\end{itemize}
\end{definition}

\begin{definition}[Principal Minors]
	Given an $n \times n$ matrix $A$, its principal minor of order $k$ ($PM_k$) is the determinant of any $k \times k$ submatrix obtained by removing $n-k$ rows and the corresponding $n-k$ columns of $A$ of the same indices.
\end{definition}

\begin{theorem}
	Let $A$ be an $n \times n$ square matrix (not necessarily symmetric). For each $k \in \{1, 2, \dots, n\}$, let $S_k$ be the sum of all principal minors of $A$ of order $k$. By convention, $S_0 = 1$. Then the characteristic polynomial of $A$ is given by:
	\[
		p_A(x) = (-1)^n \sum_{i=0}^n (-1)^i S_i x^{n-i} = \pm \left(x^n - S_1 x^{n-1} + S_2 x^{n-2} - \dots + (-1)^n S_n\right).
	\]
\end{theorem}

\begin{pfbox}
We prove the formula by induction on $n$.

Let $A\in M_{n\times n}(\mathbb{C})$. Denote by $A_{n-1}$ the upper-left
$(n-1)\times(n-1)$ submatrix of $A$. Also, for $0\le i\le n$, let
\[
S_i(A)
=
\text{the sum of all principal minors of } A \text{ of order } i,
\]
with the convention $S_0(A)=1$.

We want to show that
\[
\det(A-xI_n)
=
(-1)^n
\sum_{i=0}^n
(-1)^i S_i(A)x^{n-i}.
\]

The case $n=1$ is clear. Now assume the formula holds for $(n-1)\times(n-1)$
matrices. We prove it for an $n\times n$ matrix $A$.

Write
\[
A-xI_n=M_1+M_2,
\]
where we split the last column as
\[
\begin{pmatrix}
a_{1n}\\
\vdots\\
a_{n-1,n}\\
a_{nn}-x
\end{pmatrix}
=
\begin{pmatrix}
0\\
\vdots\\
0\\
-x
\end{pmatrix}
+
\begin{pmatrix}
a_{1n}\\
\vdots\\
a_{n-1,n}\\
a_{nn}
\end{pmatrix}.
\]
By linearity of determinant in the last column,
\[
\det(A-xI_n)=\det(M_1)+\det(M_2).
\]

First, since the last column of $M_1$ has only one nonzero entry $-x$, expanding
along the last column gives
\[
\det(M_1)
=
-x\det(A_{n-1}-xI_{n-1}).
\]
By the induction hypothesis,
\begin{align*}
\det(M_1)
&=
-x
\left[
(-1)^{n-1}
\sum_{i=0}^{n-1}
(-1)^i S_i(A_{n-1})x^{n-1-i}
\right] \\
&=
(-1)^n
\sum_{i=0}^{n-1}
(-1)^i S_i(A_{n-1})x^{n-i}.
\end{align*}

Next, compute $\det(M_2)$. When expanding $\det(M_2)$, choosing $n-i$ copies
of $-x$ from the first $n-1$ diagonal entries leaves a principal minor of order
$i$ that must contain the $n$-th row and the $n$-th column.

Let $T_i^{(n)}$ be the sum of all principal minors of $A$ of order $i$ containing the $n$-th row and column. Then
\[
\det(M_2)
=
\sum_{i=1}^n
(-1)^{n-i}T_i^{(n)}x^{n-i}.
\]

Therefore,
\begin{align*}
\det(A-xI_n)
&=
\det(M_1)+\det(M_2) \\
&=
(-1)^n
\sum_{i=0}^{n-1}
(-1)^i S_i(A_{n-1})x^{n-i}
+
\sum_{i=1}^n
(-1)^{n-i}T_i^{(n)}x^{n-i}.
\end{align*}

Now observe that every principal minor of $A$ of order $i$ either does not use
the $n$-th row and column, or does use the $n$-th row and column. Hence
\[
S_i(A)=S_i(A_{n-1})+T_i^{(n)}
\qquad
\text{for } 1\le i\le n-1.
\]
Also,
\[
S_0(A)=1,
\qquad
S_n(A)=T_n^{(n)}=\det(A).
\]
Thus we may combine the two sums and obtain
\[
\det(A-xI_n)
=
\sum_{i=0}^n
(-1)^{n-i}S_i(A)x^{n-i}.
\]
Equivalently,
\[
\det(A-xI_n)
=
(-1)^n
\sum_{i=0}^n
(-1)^i S_i(A)x^{n-i}.
\]

This proves the formula.
\end{pfbox}



\begin{theorem}[Sylvester's Criterion I]
	Let $A$ be an $n \times n$ symmetric matrix.
	\begin{enumerate}
		\item $A$ is PD if and only if $LPM_k > 0$ for every $k \in \{1, 2, \dots, n\}$.
		\item $A$ is ND if and only if $(-1)^k \cdot LPM_k > 0$ for every $k \in \{1, 2, \dots, n\}$.
	\end{enumerate}
\label{Syl. criterion I}
\end{theorem}

\begin{pfbox}
	We prove the positive-definite statement by induction on $n$.

	$(\Rightarrow)$ Assume that the $n \times n$ symmetric matrix $A$ is positive-definite. We show all its leading principal minors are positive.
	\begin{itemize}
		\item \textbf{Base case ($n=1$):} $A = (\alpha_{11})$. Since $A$ is PD, its sole eigenvalue is $\alpha_{11} > 0$. Thus $LPM_1 = \alpha_{11} > 0$, which is trivial.
		\item \textbf{Inductive step:} Suppose the statement holds for any $k \times k$ symmetric matrix. Consider $A$ to be a $(k+1) \times (k+1)$ symmetric matrix. Write $A$ as a block matrix:
		      \[
			      A = \begin{pmatrix} A_k & \mathbf{y} \\ \mathbf{y}^T & \alpha_{k+1,k+1} \end{pmatrix}
		      \]
		      where $A_k$ is the $k \times k$ upper-left submatrix. Since $A$ is PD, we have $\mathbf{x}^{T}A\mathbf{x} > 0$ for any non-zero $\mathbf{x} \in \mathbb{R}^{k+1}$. In particular, choose $\mathbf{x} = \begin{pmatrix} \mathbf{x}_k \\ 0 \end{pmatrix}$ where $\mathbf{x}_k \in \mathbb{R}^k \setminus \{\mathbf{0}\}$. Then:
		      \[
			      \mathbf{x}^T A \mathbf{x} = \begin{pmatrix} \mathbf{x}_k^T & 0 \end{pmatrix} \begin{pmatrix} A_k & \mathbf{y} \\ \mathbf{y}^T & \alpha_{k+1,k+1} \end{pmatrix} \begin{pmatrix} \mathbf{x}_k \\ 0 \end{pmatrix} = \mathbf{x}_k^T A_k \mathbf{x}_k > 0.
		      \]
		      This implies that $A_k$ is also a positive-definite matrix. By the inductive hypothesis, all leading principal minors of $A_k$ are strictly positive, which accounts for $LPM_1, LPM_2, \dots, LPM_k$. The final leading principal minor is $LPM_{k+1} = \det(A)$. Since $A$ is PD, all its eigenvalues are strictly positive, so $\det(A) = \displaystyle\prod_{i=1}^{k+1} \lambda_i > 0$. Hence, all leading principal minors of $A$ are strictly positive.
	\end{itemize}

	$(\Leftarrow)$ Conversely, suppose all the leading principal minors of an $n \times n$ symmetric matrix $A$ are positive. We show that $A$ is positive-definite by induction on $n$.
	\begin{itemize}
		\item \textbf{Base case ($n=1$):} $A = (\alpha_{11})$ with $LPM_1 = \alpha_{11} > 0$. Its eigenvalue is positive, so it is PD.
		\item \textbf{Inductive step:} Suppose this is true for any $k \times k$ symmetric matrix. Consider $A$ to be a $(k+1) \times (k+1)$ symmetric matrix whose LPMs are all strictly positive. Write:
		      \[
			      A = \begin{pmatrix} A_k & \vec{a} \\ \vec{a}^T & \beta \end{pmatrix}
		      \]
		      where $A_k \in M_k(\mathbb{R})$ is symmetric, $\vec{a} \in \mathbb{R}^k$, and $\beta \in \mathbb{R}$. Since all LPMs of $A$ are positive, all LPMs of $A_k$ are also positive. By the induction hypothesis, $A_k$ is positive-definite. Let $\{\vec{c}_1, \vec{c}_2, \dots, \vec{c}_k\}$ be the columns of $A_k$. Since $A_k$ is PD, $\det(A_k) \ne 0$, meaning these columns form a basis of $\mathbb{R}^k$. Thus, we can express $\vec{a}$ as a linear combination of the columns of $A_k$: $\vec{a} = \alpha_1 \vec{c}_1 + \dots + \alpha_k \vec{c}_k = A_k \vec{\alpha}$ where $\vec{\alpha} = (\alpha_1, \dots, \alpha_k)^T$.\\
		      We construct an elementary block matrix to eliminate $\vec{a}$ and $\vec{a}^T$:
		      \[
			      P = \begin{pmatrix} I_k & -\vec{\alpha} \\ \vec{0}^T & 1 \end{pmatrix}.
		      \]
		      Then compute the congruent transformation:
		      \[
			      P^T A P = \begin{pmatrix} I_k & \vec{0} \\ -\vec{\alpha}^T & 1 \end{pmatrix} \begin{pmatrix} A_k & \vec{a} \\ \vec{a}^T & \beta \end{pmatrix} \begin{pmatrix} I_k & -\vec{\alpha} \\ \vec{0}^T & 1 \end{pmatrix} = \begin{pmatrix} A_k & \vec{0} \\ \vec{0}^T & d \end{pmatrix}
		      \]
		      where $d = \beta - \vec{\alpha}^T A_k \vec{\alpha} \in \mathbb{R}$. Taking determinants on both sides:
		      \begin{align*}
			      \det(P^TAP)
			       & = \det(P^T)\det(A)\det(P) \\
			       & = \det(A).
		      \end{align*}
		      On the other hand, from the block diagonal form, $\det(P^T A P) = \det(A_k) \cdot d$. Therefore, $\det(A) = \det(A_k) \cdot d \implies d = \frac{\det(A)}{\det(A_k)} = \frac{LPM_{k+1}}{LPM_k}$. Since both $LPM_{k+1} > 0$ and $LPM_k > 0$, we have $d > 0$.\\
		      The eigenvalues of the block diagonal matrix $P^T A P$ are given by $\Spec(P^T A P) = \Spec(A_k) \cup \{d\}$. Since $A_k$ is PD, all its eigenvalues are positive, and since $d > 0$, all eigenvalues of $P^T A P$ are positive. Thus, $P^T A P$ is positive-definite. Since $P$ is invertible, this implies that $A$ itself must be positive-definite.
	\end{itemize}
	For the negative-definite case, note that $A$ is ND if and only if $-A$ is PD. Applying the above result to $-A$, we require $LPM_k(-A) > 0$ for all $k$. Since $-A_k$ scales every row by $-1$, $\det(-A_k) = (-1)^k \det(A_k) = (-1)^k LPM_k(A)$. Thus $(-1)^k \cdot LPM_k(A) > 0$ for all $k$.
\end{pfbox}


\begin{theorem}[Sylvester's Criterion II]
	Let $A$ be a symmetric matrix.
	\begin{enumerate}
		\item $A$ is PSD if and only if for every $k$ and every principal minor $PM_k$ of order $k$, $PM_k \ge 0$.
		\item $A$ is NSD if and only if for every $k$ and every principal minor $PM_k$ of order $k$, $(-1)^k \cdot PM_k \ge 0$.
		\item $A$ is indefinite if and only if the principal minors of $A$ do not fit into the conditions in either (1) or (2).
	\end{enumerate}
\end{theorem}


\begin{pfbox}
	$(\Rightarrow)$ The argument is similar to the proof in theorem~\eqref{Syl. criterion I}.\\
	$(\Leftarrow)$ Conversely, suppose $A$ is a symmetric matrix whose principal minors of any order are non-negative. Our goal is to prove that $A$ is PSD.

	\textbf{Claim:} For any $t > 0$, the matrix $A + tI$ is positive-definite.\\
	\textit{Proof of claim:} Consider the leading principal minors of $A+tI$. Let $A_k$ be the $k \times k$ upper-left submatrix of $A$. Then:
	$$
		LPM_k(A + tI) = \det(A_k + tI) = p_{A_k}(-t).
	$$
	Using the characteristic polynomial formula via principal minors for $A_k$:
	\[
		p_{A_k}(-t) = (-1)^k \sum_{i=0}^k (-1)^i S_i(A_k) (-t)^{k-i} = \sum_{i=0}^k S_i(A_k) t^{k-i}
	\]
	where $S_i(A_k)$ is the sum of all principal minors of order $i$ of $A_k$. Since every principal minor of $A_k$ is also a principal minor of $A$, and all principal minors of $A$ are assumed to be non-negative, we have $S_i(A_k) \ge 0$ for all $i \ge 1$. For $i = 0$, $S_0(A_k) = 1$. Since $t > 0$, every term $S_i(A_k) t^{k-i}$ is non-negative, and the $i=0$ term is $t^k > 0$. Thus, the total sum is strictly positive:
	$$
		LPM_k(A + tI) > 0 \quad \text{for every } k \in \{1, \dots, n\}.
	$$
	By Sylvester's Criterion I, $A + tI$ must be positive-definite for any $t > 0$.

	Since $A + tI$ is PD, for any non-zero vector $\mathbf{u} \in \mathbb{R}^n$, we have:
	$$
		\mathbf{u}^T (A + tI) \mathbf{u} > 0 \implies \mathbf{u}^T A \mathbf{u} + t \|\mathbf{u}\|^2 > 0.
	$$
	Taking the limit as $t \to 0^+$, we obtain:
	$$
		\mathbf{u}^T A \mathbf{u} \ge 0
	$$
	for any $\mathbf{u} \in \mathbb{R}^n$. Therefore, $A$ is positive-semidefinite.
\end{pfbox}


\begin{proposition}[General Factorization of PSD Matrices]
Let $B\in \operatorname{Sym}_n(\mathbb{R})$ and let $m\in\mathbb{N}$. Then there exists
$A\in M_{m\times n}(\mathbb{R})$ such that $B=A^tA$ if and only if $B$ is PSD
and $\rank(B)\le m$.
\end{proposition}

\begin{pfbox}
$(\Rightarrow)$ Suppose $B=A^tA$ for some $A\in M_{m\times n}(\mathbb{R})$.
For any $\mathbf{v}\in\mathbb{R}^n$,
\[
    \mathbf{v}^tB\mathbf{v}
    =
    \mathbf{v}^tA^tA\mathbf{v}
    =
    \norm{A\mathbf{v}}^2
    \ge 0.
\]
Hence $B$ is PSD. Moreover, $\ker(A^tA)=\ker(A)$.\\
By rank-nullity, $\rank(A^tA)=\rank(A)\le m$. Thus
$\rank(B)\le m$.

\medskip

$(\Leftarrow)$ Suppose $B$ is PSD and $\rank(B)=k\le m$.

If $k=0$, then $B=\mathbf{O}$, and we may choose $A=\mathbf{O}_{m\times n}$.
So assume $k\ge 1$.

By the Spectral Theorem, there exist orthonormal eigenvectors
$\mathbf{p}_1,\ldots,\mathbf{p}_k$ corresponding to the positive eigenvalues
$\lambda_1,\ldots,\lambda_k>0$ such that
\[
    B
    =
    \sum_{i=1}^k \lambda_i\mathbf{p}_i\mathbf{p}_i^t.
\]
Let $P_k=(\mathbf{p}_1\cdots\mathbf{p}_k)$ and
$\Sigma=\operatorname{diag}(\sqrt{\lambda_1},\ldots,\sqrt{\lambda_k})$. Then
$B=P_k\Sigma^2P_k^t$.

Since $m\ge k$, define
\[
    A
    =
    \left(
    \begin{array}{c}
        \Sigma P_k^t \\
        \hline
        \mathbf{O}_{(m-k)\times n}
    \end{array}
    \right)
    \in M_{m\times n}(\mathbb{R}).
\]
Then
\[
    A^tA
    =
    \left(
    \begin{array}{c|c}
        P_k\Sigma & \mathbf{O}_{n\times (m-k)}
    \end{array}
    \right)
    \left(
    \begin{array}{c}
        \Sigma P_k^t \\
        \hline
        \mathbf{O}_{(m-k)\times n}
    \end{array}
    \right)
    =
    P_k\Sigma^2P_k^t
    =
    B.
\]
Thus there exists $A\in M_{m\times n}(\mathbb{R})$ such that $B=A^tA$.
\end{pfbox}


\begin{corollary}
Let $B\in \operatorname{Sym}_n(\mathbb{R})$ be PSD and let $\rank(B)=k$.
Then there exists $A\in M_{k\times n}(\mathbb{R})$ such that
\[
    B=A^tA.
\]
Moreover, $k$ is the smallest possible number of rows of such an $A$.
\end{corollary}

\begin{pfbox}
Existence follows from the previous proposition by taking $m=k$.

For minimality, suppose $B=A^tA$ for some $A\in M_{m\times n}(\mathbb{R})$.
Then $\rank(B)=\rank(A^tA)=\rank(A)\le m$. Hence $m\ge \rank(B)=k$.
Therefore $k$ is the smallest possible number of rows.
\end{pfbox}


\begin{corollary}
Let $B\in \operatorname{Sym}_n(\mathbb{R})$. Then
\[
    B\text{ is PSD}
    \quad\Longleftrightarrow\quad
    B=A^tA
    \quad\text{for some }A\in M_n(\mathbb{R}).
\]
\end{corollary}

\begin{pfbox}
This follows from the general factorization by taking $m=n$. Indeed, if $B$ is
PSD, then $\rank(B)\le n$, so there exists $A\in M_n(\mathbb{R})$ such that
$B=A^tA$.

Conversely, if $B=A^tA$, then for any $\mathbf{v}\in\mathbb{R}^n$,
\[
    \mathbf{v}^tB\mathbf{v}
    =
    \mathbf{v}^tA^tA\mathbf{v}
    =
    \norm{A\mathbf{v}}^2
    \ge 0.
\]
Hence $B$ is PSD.
\end{pfbox}


\begin{corollary}[Symmetric Square Root]
Let $B\in \operatorname{Sym}_n(\mathbb{R})$. Then
\[
    B\text{ is PSD}
    \quad\Longleftrightarrow\quad
    B=X^2
    \quad\text{for some }X\in \operatorname{Sym}_n(\mathbb{R}).
\]
\end{corollary}

\begin{pfbox}
$(\Rightarrow)$ Suppose $B$ is PSD. By the Spectral Theorem, there exists
$P\in O(n)$ such that
\[
    B
    =
    P
    \begin{pmatrix}
        \lambda_1 &        & \mathbf{O} \\
                  & \ddots &             \\
        \mathbf{O}&        & \lambda_n
    \end{pmatrix}
    P^t,
\]
where $\lambda_i\ge 0$ for all $i$. Let
\[
    D
    =
    \begin{pmatrix}
        \lambda_1 &        & \mathbf{O} \\
                  & \ddots &             \\
        \mathbf{O}&        & \lambda_n
    \end{pmatrix},
    \qquad
    \sqrt{D}
    =
    \begin{pmatrix}
        \sqrt{\lambda_1} &        & \mathbf{O} \\
                         & \ddots &             \\
        \mathbf{O}       &        & \sqrt{\lambda_n}
    \end{pmatrix}.
\]
Choose $X=P\sqrt{D}P^t$. Then $X\in \operatorname{Sym}_n(\mathbb{R})$, and
\[
    X^2
    =
    (P\sqrt{D}P^t)(P\sqrt{D}P^t)
    =
    PD P^t
    =
    B.
\]

\medskip

$(\Leftarrow)$ Suppose $B=X^2$ for some $X\in \operatorname{Sym}_n(\mathbb{R})$.
Since $X^t=X$, we have $B=X^2=X^tX$. Hence $B$ is PSD by the previous
corollary.
\end{pfbox}


\begin{remark}
The most general factorization is the following: if $B\in\operatorname{Sym}_n(\mathbb{R})$
is PSD and $\rank(B)=k$, then $B=A^tA$ for some
$A\in M_{m\times n}(\mathbb{R})$ if and only if $m\ge k$.

The minimal choice is $m=k$. For any $m>k$, we can obtain a factorization in
$M_{m\times n}(\mathbb{R})$ by adding zero rows.

The symmetric square root is a special square factorization. If $B$ is PSD, then
$B=X^2$ for some symmetric matrix $X$. Since $X^t=X$, this also gives
$B=X^tX$ by taking $A=X$.
\end{remark}


\subsubsection*{Geometry of Quadratic Forms}
Let $A=\begin{pmatrix}a&b\\ b&c\end{pmatrix}$ be symmetric and set
$q_A(\mathbf{v})=\mathbf{v}^tA\mathbf{v}$. We study the level set
\[
    \Sigma_1(A)=\{\mathbf{v}\in\mathbb{R}^2:q_A(\mathbf{v})=1\}.
\]
The cross term $2bxy$ is geometrically a rotation of coordinates. 
By Spectral Theorem, there is an ONB $\{\mathbf{u}_1,\mathbf{u}_2\}$ such that, if
$\mathbf{v}=s\mathbf{u}_1+t\mathbf{u}_2$, then
\[
    q_A(\mathbf{v})=\lambda_1s^2+\lambda_2t^2.
\]
Thus the shape of $\Sigma_1(A)$ is determined by the signs of
$\lambda_1,\lambda_2$.

\begin{enumerate}
    \item \textbf{PD.} If $\lambda_1,\lambda_2>0$, then
    $\lambda_1s^2+\lambda_2t^2=1$ is an ellipse. Its semi-axis lengths are
    $1/\sqrt{\lambda_1}$ and $1/\sqrt{\lambda_2}$, so the longer axis corresponds
    to the smaller eigenvalue.

    \begin{center}
    \begin{tikzpicture}[scale=1, every node/.style={font=\small}]
        \draw[->,gray75] (-2.2,0)--(2.2,0) node[right] {$x$};
        \draw[->,gray75] (0,-2)--(0,2) node[above] {$y$};
        \draw[greenblue,thick,rotate=62] (0,0) ellipse (1.35 and 0.66);
        \draw[->,greenblue!70!black,thick]
            (0,0)--({1.35*cos(62)},{1.35*sin(62)})
            node[above right] {$\mathbf{v}_1$};
        \draw[->,greenblue!70!black,thick]
            (0,0)--({-0.66*sin(62)},{0.66*cos(62)})
            node[above left] {$\mathbf{v}_2$};
    \end{tikzpicture}
    \end{center}

    For example, take
    $A=\begin{pmatrix}4&-\frac{3}{2}\\ -\frac{3}{2}&2\end{pmatrix}$. Then
    $\Sigma_1(A)$ is $4x^2-3xy+2y^2=1$. Since the eigenvalues of $A$ are
    $(6+\sqrt{13})/2$ and $(6-\sqrt{13})/2$, this is an ellipse. The cross term
    $-3xy$ rotates the axes of the ellipse.

    \item \textbf{Indefinite.} If $\lambda_1>0>\lambda_2$, then
    $\lambda_1s^2+\lambda_2t^2=1$ is a hyperbola. The positive eigendirection is
    the direction in which the level $1$ opens. The equation $q_A(\mathbf{v})=0$
    gives the two asymptotic directions.

    \begin{center}
    \begin{tikzpicture}[scale=0.95, every node/.style={font=\small}]
        \draw[->,gray75] (-2.8,0)--(2.8,0) node[right] {$x$};
        \draw[->,gray75] (0,-2.2)--(0,2.2) node[above] {$y$};
        \begin{scope}[rotate=45]
            \draw[gray75,densely dashed,thick] (-1.15,-1.99)--(1.15,1.99)
                node[pos=0.88,above] {asymptote};
            \draw[gray75,densely dashed,thick] (-1.15,1.99)--(1.15,-1.99)
                node[pos=0.87,below] {asymptote};
            \draw[darkred,thick,domain=-2.05:2.05,smooth,variable=\x]
                plot ({sqrt(1+\x*\x)/sqrt(3)},{\x});
            \draw[darkred,thick,domain=-2.05:2.05,smooth,variable=\x]
                plot ({-sqrt(1+\x*\x)/sqrt(3)},{\x});
            \draw[->,greenblue!70!black,thick] (0,0)--(1.45,0)
                node[below right] {$\mathbf{v}_+$};
            \draw[->,mountbattenpink,thick] (0,0)--(0,1.35)
                node[left] {$\mathbf{v}_-$};
        \end{scope}
    \end{tikzpicture}
    \end{center}

    For example, if $A=\begin{pmatrix}1&2\\2&1\end{pmatrix}$, then
    $\Sigma_1(A)$ is $x^2+4xy+y^2=1$. In the eigenbasis
    $\mathbf{v}_+=(1,1)/\sqrt{2}$ and $\mathbf{v}_-=(1,-1)/\sqrt{2}$, this becomes
    $3s^2-t^2=1$. The dashed lines come from $3s^2-t^2=0$, so they are the
    asymptotes.

    Think first about directions. For a direction
    $\mathbf{w}=s\mathbf{v}_+ + t\mathbf{v}_-$, the sign of
    $q_A(\mathbf{w})=3s^2-t^2$ tells us whether this direction can hit the level
    $q_A=1$ after scaling. If $3s^2-t^2>0$, the $\mathbf{v}_+$ direction accounts
    more, and $c\mathbf{w}$ can be chosen so that $q_A(c\mathbf{w})=1$. If
    $3s^2-t^2\le0$, then no positive scaling can make $q_A=1$. Near an asymptote,
    $3s^2-t^2$ is positive but very small, so the needed scaling
    $c=1/\sqrt{q_A(\mathbf{w})}$ is very large; this is why the branch stretches
    far out. Near the opening direction $\mathbf{v}_+$, the positive part clearly
    dominates, so only a small scaling is needed, and the branch stays close to
    the vertex/focus region.

    \item \textbf{PSD.} If $\lambda_1>0$ and $\lambda_2=0$, then
    $\lambda_1s^2=1$, so $s=\pm 1/\sqrt{\lambda_1}$ and $t$ is free. Therefore
    $\Sigma_1(A)$ is two parallel lines, both parallel to the null direction
    $\ker(A)$.

    \begin{center}
    \begin{tikzpicture}[scale=0.95, every node/.style={font=\small}]
        \draw[->,gray75] (-2.2,0)--(2.2,0) node[right] {$x$};
        \draw[->,gray75] (0,-2.1)--(0,2.1) node[above] {$y$};
        \draw[asparagus,thick] (-1.05,2.05)--(2.05,-1.05);
        \draw[asparagus,thick] (-2.05,1.05)--(1.05,-2.05);
        \draw[->,asparagus!65!black,thick] (0,0)--(1.25,-1.25)
            node[right] {$\ker(A)$};
        \node[above right] at (1.25,-0.25) {$x+y=1$};
        \node[below left] at (-1.25,0.25) {$x+y=-1$};
    \end{tikzpicture}
    \end{center}

    For example, if $A=\begin{pmatrix}1&1\\1&1\end{pmatrix}$, then
    $q_A(x,y)=(x+y)^2$, so $\Sigma_1(A)$ is $(x+y)^2=1$, i.e.
    $x+y=1$ and $x+y=-1$. These lines are parallel to
    $\ker(A)=\operatorname{span}\{(1,-1)\}$.

    \item \textbf{ND/NSD.} If $\lambda_1,\lambda_2\le 0$, then
    $q_A(\mathbf{v})\le 0$ for every $\mathbf{v}\in\mathbb{R}^2$. Hence the
    positive level set $\Sigma_1(A)$ is empty.
\end{enumerate}

\subsubsection*{Application: Second Derivative Tests}

\begin{theorem}
	Let $f(x_1, x_2, \dots, x_n)$ be a function with continuous third-order derivatives, and let $\mathbf{a} = (a_1, a_2, \dots, a_n)$ be a critical point of $f$, i.e., $\frac{\partial f}{\partial x_i}(\mathbf{a}) = 0$ for each $i = 1, 2, \dots, n$. Consider the $n \times n$ Hessian matrix at $\mathbf{a}$:
	$$
		H_f(\mathbf{a}) = \left( \frac{\partial^2 f}{\partial x_i \partial x_j}(\mathbf{a}) \right)_{i,j}.
	$$
	\begin{enumerate}
		\item If $H_f(\mathbf{a})$ is PD, then $\mathbf{a}$ is a local minimum point for $f$.
		\item If $H_f(\mathbf{a})$ is ND, then $\mathbf{a}$ is a local maximum point for $f$.
		\item If $H_f(\mathbf{a})$ is indefinite (has at least one strictly positive and one strictly negative eigenvalue), then $\mathbf{a}$ is a saddle point.
		\item In any other scenario (i.e., when $H_f(\mathbf{a})$ is PSD or NSD but not definite), the test is inconclusive.
	\end{enumerate}
\end{theorem}

\begin{pfbox}
	We prove statement (1). By shifting the coordinate system and the function value suitably, we can assume without loss of generality that $\mathbf{a} = \mathbf{0}$ and $f(\mathbf{0}) = 0$. By multivariable Taylor's Theorem, writing $\mathbf{x} = (x_1, x_2, \dots, x_n)$, we expand around $\mathbf{0}$:
	\[
		f(\mathbf{x}) = f(\mathbf{0}) + \sum_{i=1}^n \frac{\partial f}{\partial x_i}(\mathbf{0})x_i + \frac{1}{2}\sum_{i,j=1}^n \frac{\partial^2 f}{\partial x_i \partial x_j}(\mathbf{0})x_i x_j + S(\mathbf{x})
	\]
	where $S(\mathbf{x})$ is the remainder term satisfying $\lim_{\mathbf{x}\to\mathbf{0}} \frac{S(\mathbf{x})}{\|\mathbf{x}\|^2} = 0$. Since $\mathbf{0}$ is a critical point, the first-order sum vanishes. Thus:
	$$
		f(\mathbf{x}) = \frac{1}{2} \mathbf{x}^T H_f(\mathbf{0}) \mathbf{x} + S(\mathbf{x}).
	$$
	By the Real Spectral Theorem, since $H_f(\mathbf{0})$ is symmetric, there exists an ONB $\{\mathbf{v}_1, \dots, \mathbf{v}_n\}$ of $\mathbb{R}^n$ consisting of eigenvectors of $H_f(\mathbf{0})$. We can express the quadratic form in terms of the eigenvalues $\lambda_i \in Spec(H_f(\mathbf{0}))$:
	$$
		f(\mathbf{x}) = \frac{1}{2}\sum_{i=1}^n \lambda_i \innerproduct{\mathbf{x}}{\mathbf{v}_i}^2 + S(\mathbf{x}).
	$$
	Assume $H_f(\mathbf{0})$ is positive-definite. Then all eigenvalues satisfy $\lambda_i > 0$. Take any $\varepsilon > 0$ chosen small enough such that $0 < \varepsilon < \frac{\lambda_i}{2}$ for all $i \in \{1, \dots, n\}$.\\
	Since $\lim_{\mathbf{x}\to\mathbf{0}} \frac{S(\mathbf{x})}{\|\mathbf{x}\|^2} = 0$, there exists a $\delta > 0$ such that:
	\begin{align*}
		\|\mathbf{x}\|<\delta
		 & \implies
		\left|\frac{S(\mathbf{x})}{\|\mathbf{x}\|^2}\right|<\varepsilon \\
		 & \implies
		|S(\mathbf{x})|<\varepsilon\|\mathbf{x}\|^2.
	\end{align*}
	This implies $S(\mathbf{x}) > -\varepsilon \|\mathbf{x}\|^2$. Substituting this into our expression for $f(\mathbf{x})$ yields:
	$$
		f(\mathbf{x}) > \frac{1}{2}\sum_{i=1}^n \lambda_i \innerproduct{\mathbf{x}}{\mathbf{v}_i}^2 - \varepsilon \|\mathbf{x}\|^2.
	$$
	Let us write $\mathbf{x} = \displaystyle\sum_{i=1}^n \alpha_i \mathbf{v}_i$, where $\alpha_i = \innerproduct{\mathbf{x}}{\mathbf{v}_i}$. By Pythagorean theorem, $\|\mathbf{x}\|^2 = \displaystyle\sum_{i=1}^n \alpha_i^2$. Substituting these representations gives:
	\begin{align*}
		f(\mathbf{x})
		 & > \frac{1}{2}\sum_{i=1}^n \lambda_i\alpha_i^2
		- \varepsilon\sum_{i=1}^n \alpha_i^2                                       \\
		 & = \sum_{i=1}^n \left(\frac{1}{2}\lambda_i-\varepsilon\right)\alpha_i^2.
	\end{align*}
	Since we picked $\varepsilon < \frac{\lambda_i}{2}$, the coefficient $\left(\frac{1}{2}\lambda_i - \varepsilon\right)$ is strictly positive for every $i$. Therefore, for any $\mathbf{x} \ne \mathbf{0}$ within the neighborhood $\|\mathbf{x}\| < \delta$, at least one $\alpha_i \ne 0$, which ensures:
	$$
		f(\mathbf{x}) > 0 = f(\mathbf{0}).
	$$
	This proves that $\mathbf{a} = \mathbf{0}$ is a strict local minimum point for $f$.
\end{pfbox}

\begin{example}[2-Variable Function]
	Let $f(x,y)$ be a $C^3$-function and fix $(a,b)$ in the domain of $f$. Then one can write
	\begin{align*}
		f(x,y) & = f(a,b) + f_x(a,b)(x-a) + f_y(a,b)(y-b)                                                                  \\
		       & \quad + \frac{1}{2}\left(f_{xx}(a,b)(x-a)^2 + 2f_{xy}(a,b)(x-a)(y-b) + f_{yy}(a,b)(y-b)^2\right) + R(x,y)
	\end{align*}
	where $R(x,y)$ is a remainder term that satisfies
	$$
		\lim_{(x,y)\to(a,b)} \frac{R(x,y)}{(x-a)^2 + (y-b)^2} = 0.
	$$
\end{example}

Suppose $(0,0)$ is a critical point, i.e., $f_x(0,0) = 0$ and $f_y(0,0) = 0$. Then the local behavior can be approximated by the quadratic form:
$$
	f(x,y) - f(0,0) \approx \frac{1}{2}\left(Ax^2 + 2Bxy + Cy^2\right)
$$
where $A = f_{xx}(0,0)$, $B = f_{xy}(0,0)$, and $C = f_{yy}(0,0)$. This quadratic form is associated with the Hessian matrix:
\[
	\begin{pmatrix} A & B \\ B & C \end{pmatrix}.
\]
The signs of the eigenvalues $\lambda_1, \lambda_2$ of the Hessian matrix determine the nature of the critical point:
\begin{itemize}
	\item If $\lambda_1, \lambda_2 > 0$, then $f(x,y) > f(0,0)$ near $(0,0)$, so $(0,0)$ is a local minimum.
	\item If $\lambda_1, \lambda_2 < 0$, then $f(x,y) < f(0,0)$ near $(0,0)$, so $(0,0)$ is a local maximum.
	\item If $\lambda_1 \lambda_2 < 0$, then $(0,0)$ is a saddle point.
\end{itemize}

\subsubsection*{Convexity and Global Minimum}

The second derivative test is local: it tells us what happens near a critical point. Convexity gives a global version of this idea.

\begin{definition}[Convex Set]
	A subset $U \subseteq \mathbb{R}^n$ is called \textbf{convex} if for every $\mathbf{x}, \mathbf{y} \in U$ and every $t \in [0, 1]$, we have:
	$$
		(1-t)\mathbf{x} + t\mathbf{y} \in U.
	$$
	Geometrically, this means that the straight line segment joining any two points of $U$ is entirely contained within $U$.
\end{definition}

\begin{definition}[Convex Function]
	Let $U \subseteq \mathbb{R}^n$ be a convex set.
	\begin{enumerate}
		\item A function $f: U \to \mathbb{R}$ is called \textbf{convex} if for every $\mathbf{x}, \mathbf{y} \in U$ and every $t \in [0,1]$:
		      $$
			      f((1-t)\mathbf{x} + t\mathbf{y}) \le (1-t)f(\mathbf{x}) + tf(\mathbf{y}).
		      $$
		\item If the inequality is strict whenever $\mathbf{x} \ne \mathbf{y}$ and $0 < t < 1$, then $f$ is called \textbf{strictly convex}.
	\end{enumerate}
\end{definition}

\begin{theorem}[Second-Order Criterion for Convexity]
	Let $U \subseteq \mathbb{R}^n$ be open and convex, and let $f: U \to \mathbb{R}$ be a $C^2$-function. Then:
	\begin{enumerate}
		\item $f$ is convex if and only if $H_f(\mathbf{x})$ is PSD for all $\mathbf{x} \in U$.
		\item If $H_f(\mathbf{x})$ is positive-definite for every $\mathbf{x} \in U$, then $f$ is strictly convex.
	\end{enumerate}
\end{theorem}

\begin{pfbox}
	We can connect multivariable convexity to single-variable calculus using line restrictions. Let $\mathbf{a} \in U$ and choose a direction vector $\mathbf{b} \in \mathbb{R}^n$. Define the single-variable function $g(t) = f(\mathbf{a} + t\mathbf{b})$.

	Applying the chain rule, the first and second derivatives of $g(t)$ are:
	\begin{align*}
		g'(t)
		 & = \sum_{i=1}^n \frac{\partial f}{\partial x_i}(\mathbf{a}+t\mathbf{b})b_i \\
		 & = \nabla f(\mathbf{a}+t\mathbf{b})\cdot\mathbf{b}.
	\end{align*}
	\begin{align*}
		g''(t)
		 & = \sum_{j=1}^n\sum_{i=1}^n
		\frac{\partial^2 f}{\partial x_i\partial x_j}(\mathbf{a}+t\mathbf{b})b_i b_j \\
		 & = \mathbf{b}^T H_f(\mathbf{a}+t\mathbf{b})\mathbf{b}.
	\end{align*}

	$(\Rightarrow)$ Suppose $f$ is convex. Then the single-variable restriction $g(t)$ must be a convex function on its interval of definition. From single-variable calculus, a $C^2$ function is convex if and only if its second derivative is non-negative everywhere, meaning $g''(t) \ge 0$. Evaluating at $t=0$ gives $g''(0) = \mathbf{b}^T H_f(\mathbf{a}) \mathbf{b} \ge 0$. Since this holds for any vector $\mathbf{b} \in \mathbb{R}^n$ and any point $\mathbf{a} \in U$, the Hessian matrix $H_f(\mathbf{x})$ is positive-semidefinite for all $\mathbf{x} \in U$.

	$(\Leftarrow)$ Conversely, suppose $H_f(\mathbf{x})$ is PSD for all $\mathbf{x} \in U$. Take any two points $\mathbf{x}, \mathbf{y} \in U$. Since $U$ is a convex set, the segment $t\mathbf{x} + (1-t)\mathbf{y}$ stays in $U$ for all $t \in [0,1]$. Define $g(t) = f(\mathbf{y} + t(\mathbf{x}-\mathbf{y}))$. Its second derivative is given by:
	$$
		g''(t) = (\mathbf{x}-\mathbf{y})^T H_f(t\mathbf{x} + (1-t)\mathbf{y}) (\mathbf{x}-\mathbf{y}).
	$$
	Since $H_f$ is PSD everywhere on $U$, we know that $g''(t) \ge 0$ for all $t \in [0,1]$. By single-variable calculus, $g(t)$ is a convex function on $[0,1]$, which satisfies:
	\begin{align*}
		g(t)
		 & = g(t\cdot 1+(1-t)\cdot 0) \\
		 & \le tg(1)+(1-t)g(0).
	\end{align*}
	Substituting the definition of $g(t)$ back into the inequality:
	$$
		f(t\mathbf{x} + (1-t)\mathbf{y}) \le t f(\mathbf{x}) + (1-t) f(\mathbf{y}).
	$$
	This matches the definition of a convex function, proving that $f$ is convex on $U$.
\end{pfbox}

\begin{theorem}[Convex Functions and Global Minima]
	Let $U \subseteq \mathbb{R}^n$ be a convex set, and let $f: U \to \mathbb{R}$ be a differentiable and convex function. If $\mathbf{a} \in U$ satisfies $\nabla f(\mathbf{a}) = \mathbf{0}$, then $\mathbf{a}$ is a global minimum of $f$ on $U$.
\end{theorem}

\begin{remark}
	Furthermore, if $f$ is strictly convex, then this global minimum is unique.
\end{remark}

\begin{pfbox}
	Let $\mathbf{a}$ be a critical point of $f$, so $\nabla f(\mathbf{a}) = \mathbf{0}$. We want to show that $f(\mathbf{y}) \ge f(\mathbf{a})$ for any arbitrary point $\mathbf{y} \in U$. Let us construct the auxiliary line segment function $g(t) = f((1-t)\mathbf{a} + t\mathbf{y}) = f(\mathbf{a} + t(\mathbf{y}-\mathbf{a}))$ defined for $t \in [0, 1]$. Since $U$ is a convex set, this path lies entirely inside $U$.\\
	Because $f$ is given to be a convex function, its line restriction $g(t)$ is a convex function of a single variable on $[0, 1]$. A differentiable single-variable function is convex if and only if its derivative $g'(t)$ is a monotonically increasing function. Therefore, we have:
	$$
		g'(t) \ge g'(0) \quad \text{for all } t \in [0, 1].
	$$
	Integrating this derivative inequality from $t=0$ to $t=1$:
	$$
		\int_0^1 g'(t) \, dt \ge \int_0^1 g'(0) \, dt \implies g(1) - g(0) \ge g'(0) \cdot (1 - 0) = g'(0).
	$$
	Let us compute the value of $g'(0)$ explicitly via the multivariable chain rule:
	$$
		g'(t) = \nabla f(\mathbf{a} + t(\mathbf{y}-\mathbf{a})) \cdot (\mathbf{y}-\mathbf{a}) \implies g'(0) = \nabla f(\mathbf{a}) \cdot (\mathbf{y}-\mathbf{a}).
	$$
	Since $\mathbf{a}$ is a critical point, $\nabla f(\mathbf{a}) = \mathbf{0}$, which yields $g'(0) = \mathbf{0} \cdot (\mathbf{y}-\mathbf{a}) = 0$. Substituting this back into the integrated expression:
	$$
		g(1) - g(0) \ge 0 \implies g(1) \ge g(0).
	$$
	Since $g(1) = f(\mathbf{y})$ and $g(0) = f(\mathbf{a})$, we conclude that:
	$$
		f(\mathbf{y}) \ge f(\mathbf{a})
	$$
	for all $\mathbf{y} \in U$. This establishes that $\mathbf{a}$ is a global minimum of $f$ on $U$.
\end{pfbox}



\subsection{Min--Max Theorem \& Weyl's Inequality}


\begin{proposition}
Suppose $A\in M_n(\mathbb{C})$ is a positive-definite Hermitian matrix whose eigenvalues are given by
$\lambda_1\ge \lambda_2\ge \dots \ge \lambda_n$.
Let the Rayleigh quotient defined by
\[
f(\mathbf{x})=\frac{\mathbf{x}^* A\mathbf{x}}{\mathbf{x}^*\mathbf{x}}\quad \forall \mathbf{x}\ne \mathbf{0}.
\]
Then
\[
    \lambda_1
    =
    \max_{\substack{\mathbf{x}\in\mathbb{C}^n\\ \mathbf{x}\neq \mathbf{0}}}
    f(\mathbf{x}),
    \qquad
    \lambda_n
    =
    \min_{\substack{\mathbf{x}\in\mathbb{C}^n\\ \mathbf{x}\neq \mathbf{0}}}
    f(\mathbf{x}).
\]
\end{proposition}

\begin{pfbox}
By spectral theorem, there exists an ONB $\{\mathbf{v}_1,\dots,\mathbf{v}_n\}$ of eigenvectors of $A$. Write $\mathbf{x}=\displaystyle\sum_{i=1}^n \innerproduct{\mathbf{x}}{\mathbf{v}_i}\mathbf{v}_i = \displaystyle\sum_{i=1}^n \alpha_i\mathbf{v}_i$, then
\begin{align*}
    f(\mathbf{x})
    &=
    \frac{
    \bp{\displaystyle\sum_i \alpha_i\mathbf{v}_i}^*
    A\bp{\displaystyle\sum_j \alpha_j\mathbf{v}_j}
    }
    {\norm{\mathbf{x}}^2} \\
    &=
    \frac{
    \bp{\displaystyle\sum_i \alpha_i\mathbf{v}_i}^*
    \bp{\displaystyle\sum_j \alpha_j A\mathbf{v}_j}
    }
    {\displaystyle\sum_i \abs{\alpha_i}^2} \\
    &=
    \frac{
    \innerproduct{\displaystyle\sum_i \alpha_i\mathbf{v}_i}
    {\displaystyle\sum_j \alpha_j\lambda_j\mathbf{v}_j}
    }
    {\displaystyle\sum_i \abs{\alpha_i}^2} \\
    &=
    \frac{
    \displaystyle\sum_i\sum_j \alpha_i\overline{\alpha_j}\lambda_j
    \innerproduct{\mathbf{v}_i}{\mathbf{v}_j}
    }
    {\displaystyle\sum_i \abs{\alpha_i}^2}.
\end{align*}
Since $\{\mathbf{v}_1,\dots,\mathbf{v}_n\}$ is an ONB, only the terms with
$i=j$ survive. Hence
\[
    f(\mathbf{x})
    =
    \frac{
    \displaystyle\sum_{i=1}^n \abs{\alpha_i}^2\lambda_i
    }
    {
    \displaystyle\sum_{i=1}^n \abs{\alpha_i}^2
    }.
\]
Thus the Rayleigh quotient is a weighted average of
$\lambda_1,\lambda_2,\dots,\lambda_n$.

In particular,
\[
    \lambda_n
    \le
    f(\mathbf{x})
    \le
    \lambda_1.
\]
Moreover,
\[
    \lambda_1=f(\mathbf{v}_1),
    \qquad
    \lambda_n=f(\mathbf{v}_n).
\]
Therefore,
\[
    \max_{\mathbf{x}\ne \mathbf{0}} f(\mathbf{x})=\lambda_1,
    \qquad
    \min_{\mathbf{x}\ne \mathbf{0}} f(\mathbf{x})=\lambda_n.
\]
\end{pfbox}

\begin{theorem}[Courant--Fischer Min--Max]
Let $A$ be Hermitian on a complex inner product space
$(V,\innerproduct{-}{-})$ with $\dim V=n$. Suppose the eigenvalues of $A$
are ordered as
\[
    \lambda_1\ge \lambda_2\ge \cdots \ge \lambda_n.
\]
Then, for each $k=1,2,\dots,n$,
\[
    \lambda_k
    =
    \min_{\substack{E\subseteq V \\ \dim E=n-k+1}}
    \max_{\substack{\mathbf{x}\in E \\ \mathbf{x}\ne \mathbf{0}}}
    \frac{\innerproduct{\mathbf{x}}{A\mathbf{x}}}
         {\innerproduct{\mathbf{x}}{\mathbf{x}}}
    =
    \max_{\substack{F\subseteq V \\ \dim F=k}}
    \min_{\substack{\mathbf{x}\in F \\ \mathbf{x}\ne \mathbf{0}}}
    \frac{\innerproduct{\mathbf{x}}{A\mathbf{x}}}
         {\innerproduct{\mathbf{x}}{\mathbf{x}}}.
\]
\end{theorem}

\begin{pfbox}
For $\mathbf{x}\neq \mathbf{0}$, write
\[
    f_A(\mathbf{x})
    =
    \frac{\innerproduct{\mathbf{x}}{A\mathbf{x}}}
         {\innerproduct{\mathbf{x}}{\mathbf{x}}}.
\]
Since $f_A(c\mathbf{x})=f_A(\mathbf{x})$ for $c\neq 0$, it suffices to consider
unit vectors.

By the Spectral Theorem, there exists an ONB
$\{\mathbf{u}_1,\ldots,\mathbf{u}_n\}$ of eigenvectors of $A$ such that
$A\mathbf{u}_i=\lambda_i\mathbf{u}_i$, where
$\lambda_1\ge \lambda_2\ge \cdots \ge \lambda_n$.

\medskip

\textbf{Proof of the first equality.}

We prove
\[
    \lambda_k
    =
    \min_{\substack{E\subseteq V\\ \dim E=n-k+1}}
    \max_{\substack{\mathbf{x}\in E\\ \norm{\mathbf{x}}=1}}
    \innerproduct{\mathbf{x}}{A\mathbf{x}}.
\]

\textbf{Case 1.}  
Let $E_0=\spn\{\mathbf{u}_k,\mathbf{u}_{k+1},\ldots,\mathbf{u}_n\}$.
Then $\dim E_0=n-k+1$.

For any $\mathbf{x}\in E_0$ with $\norm{\mathbf{x}}=1$, write
$\mathbf{x}=\displaystyle\sum_{i=k}^n \alpha_i\mathbf{u}_i$. Then
$\displaystyle\sum_{i=k}^n \abs{\alpha_i}^2=1$. Hence
\[
    \innerproduct{\mathbf{x}}{A\mathbf{x}}
    =
    \sum_{i=k}^n \abs{\alpha_i}^2\lambda_i
    \le
    \sum_{i=k}^n \abs{\alpha_i}^2\lambda_k
    =
    \lambda_k.
\]
On the other hand, taking $\mathbf{x}=\mathbf{u}_k$ gives
$\innerproduct{\mathbf{u}_k}{A\mathbf{u}_k}=\lambda_k$. Therefore,
\[
    \max_{\substack{\mathbf{x}\in E_0\\ \norm{\mathbf{x}}=1}}
    \innerproduct{\mathbf{x}}{A\mathbf{x}}
    =
    \lambda_k.
\]
Thus,
\[
    \min_{\substack{E\subseteq V\\ \dim E=n-k+1}}
    \max_{\substack{\mathbf{x}\in E\\ \norm{\mathbf{x}}=1}}
    \innerproduct{\mathbf{x}}{A\mathbf{x}}
    \le
    \lambda_k.
\]

\textbf{Case 2.}  
Let $E\subseteq V$ be any subspace with $\dim E=n-k+1$.\\
Set $V_k=\spn\{\mathbf{u}_1,\ldots,\mathbf{u}_k\}$, so $\dim V_k=k$.
We claim that $V_k\cap E\neq \{\mathbf{0}\}$. 
Indeed, if $V_k\cap E=\{\mathbf{0}\}$, then
\[
    \dim(V_k+E)
    =
    \dim V_k+\dim E-\dim(V_k\cap E)
    =
    k+(n-k+1)-0
    =
    n+1,
\]
which is impossible since $V_k+E\subseteq V$. Hence, there exists $\mathbf{u}\in E$ such that $\mathbf{u}=\displaystyle\sum_{i=1}^k \beta_i\mathbf{u}_i$, where
$\displaystyle\sum_{i=1}^k\abs{\beta_i}^2=1$. Then
\[
    \innerproduct{\mathbf{u}}{A\mathbf{u}}
    =
    \sum_{i=1}^k \abs{\beta_i}^2\lambda_i
    \ge
    \sum_{i=1}^k \abs{\beta_i}^2\lambda_k
    =
    \lambda_k.
\]
So we have
\[
    \max_{\substack{\mathbf{x}\in E\\ \norm{\mathbf{x}}=1}}
    \innerproduct{\mathbf{x}}{A\mathbf{x}}
    \ge
    \lambda_k.
\]
Because $E$ was arbitrary,
\[
    \min_{\substack{E\subseteq V\\ \dim E=n-k+1}}
    \max_{\substack{\mathbf{x}\in E\\ \norm{\mathbf{x}}=1}}
    \innerproduct{\mathbf{x}}{A\mathbf{x}}
    \ge
    \lambda_k.
\]

Combining the two cases, we obtain
\[
    \lambda_k
    =
    \min_{\substack{E\subseteq V\\ \dim E=n-k+1}}
    \max_{\substack{\mathbf{x}\in E\\ \norm{\mathbf{x}}=1}}
    \innerproduct{\mathbf{x}}{A\mathbf{x}}.
\]

\medskip

\textbf{Proof of the second equality.}

Apply the first equality to $-A$. Since the eigenvalues of $-A$ are
$-\lambda_n\ge -\lambda_{n-1}\ge \cdots \ge -\lambda_1$, we have
$\lambda_{n-k+1}(-A)=-\lambda_k(A)$. Therefore,
\[
    -\lambda_k(A)
    =
    \min_{\substack{F\subseteq V\\ \dim F=k}}
    \max_{\substack{\mathbf{x}\in F\\ \norm{\mathbf{x}}=1}}
    \innerproduct{\mathbf{x}}{(-A)\mathbf{x}}.
\]
That is,
\[
    -\lambda_k(A)
    =
    \min_{\substack{F\subseteq V\\ \dim F=k}}
    \max_{\substack{\mathbf{x}\in F\\ \norm{\mathbf{x}}=1}}
    \bp{-\innerproduct{\mathbf{x}}{A\mathbf{x}}}.
\]
Taking negative signs on both sides gives
\begin{align*}
    \lambda_k(A)
    &=
    -\min_{\substack{F\subseteq V\\ \dim F=k}}
    \max_{\substack{\mathbf{x}\in F\\ \norm{\mathbf{x}}=1}}
    \bp{-\innerproduct{\mathbf{x}}{A\mathbf{x}}}                         \\
    &=
    \max_{\substack{F\subseteq V\\ \dim F=k}}
    \min_{\substack{\mathbf{x}\in F\\ \norm{\mathbf{x}}=1}}
    \innerproduct{\mathbf{x}}{A\mathbf{x}}.
\end{align*}
\end{pfbox}

\begin{remark}
Since
\[
    \frac{\innerproduct{c\mathbf{x}}{A(c\mathbf{x})}}
         {\innerproduct{c\mathbf{x}}{c\mathbf{x}}}
    =
    \frac{\innerproduct{\mathbf{x}}{A\mathbf{x}}}
         {\innerproduct{\mathbf{x}}{\mathbf{x}}}
    \qquad (c\neq 0),
\]
the Rayleigh quotient only depends on the direction of $\mathbf{x}$, not on
its length. Hence we may restrict to unit vectors.

When $k=1$, the first formula degenerates to
\[
    \lambda_1
    =
    \max_{\substack{\mathbf{x}\in V\\ \mathbf{x}\neq \mathbf{0}}}
    \frac{\innerproduct{\mathbf{x}}{A\mathbf{x}}}
         {\innerproduct{\mathbf{x}}{\mathbf{x}}},
\]
because the only $n$-dimensional subspace of $V$ is $V$ itself.

When $k=n$, the first formula degenerates to
\[
    \lambda_n
    =
    \min_{\substack{\dim E=1}}
    \max_{\substack{\mathbf{x}\in E\\ \mathbf{x}\neq \mathbf{0}}}
    \frac{\innerproduct{\mathbf{x}}{A\mathbf{x}}}
         {\innerproduct{\mathbf{x}}{\mathbf{x}}}.
\]
But a one-dimensional subspace is of the form $E=\spn\{\mathbf{v}\}$.
Since the quotient only depends on direction, the inner maximum over
$\spn\{\mathbf{v}\}$ is just the value at any nonzero vector in that line.
Thus,
\[
    \lambda_n
    =
    \min_{\substack{\mathbf{v}\in V\\ \mathbf{v}\neq \mathbf{0}}}
    \frac{\innerproduct{\mathbf{v}}{A\mathbf{v}}}
         {\innerproduct{\mathbf{v}}{\mathbf{v}}}.
\]
Similarly, the second formula gives
\[
    \lambda_1
    =
    \max_{\substack{\dim F=1}}
    \min_{\substack{\mathbf{x}\in F\\ \mathbf{x}\neq \mathbf{0}}}
    \frac{\innerproduct{\mathbf{x}}{A\mathbf{x}}}
         {\innerproduct{\mathbf{x}}{\mathbf{x}}},
\]
which is the same as maximizing over all directions.
\end{remark}


\begin{corollary}
Let $A\in \operatorname{Sym}_n(\mathbb{R})$ and define $f_A(\mathbf{x})=\mathbf{x}^tA\mathbf{x}$.
Assume the eigenvalues of $A$ are ordered as $\lambda_1\ge \lambda_2\ge \cdots \ge \lambda_n$.
Then, for $1\le k\le n$,
\[
    \lambda_k
    =
    \min_{\substack{E\subseteq \mathbb{R}^n\\ \dim E=n-k+1}}
    \max_{\substack{\mathbf{x}\in E\\ \norm{\mathbf{x}}=1}}
    f_A(\mathbf{x})
    =
    \max_{\substack{E\subseteq \mathbb{R}^n\\ \dim F=k}}
    \min_{\substack{\mathbf{x}\in E\\ \norm{\mathbf{x}}=1}}
    f_A(\mathbf{x}).
\]
\end{corollary}

\begin{proposition}
For Hermitian matrices $A,H\in M_n(\mathbb{C})$,
\[
    \lambda_1(A+H)-\lambda_1(A)\le\lambda_1(H).
\]
\label{prop:first-eigenvalue-diff-upper-bound}
\end{proposition}

\begin{pfbox}
Using the Rayleigh quotient characterization of the largest eigenvalue,
\begin{align*}
    \lambda_1(A+H)
    &=
    \max_{\norm{\mathbf{x}}=1}
    \innerproduct{\mathbf{x}}{(A+H)\mathbf{x}}                         \\
    &=
    \max_{\norm{\mathbf{x}}=1}
    \bp{
    \innerproduct{\mathbf{x}}{A\mathbf{x}}
    +
    \innerproduct{\mathbf{x}}{H\mathbf{x}}
    }                                                                  \\
    &\le
    \max_{\norm{\mathbf{x}}=1}
    \innerproduct{\mathbf{x}}{A\mathbf{x}}
    +
    \max_{\norm{\mathbf{x}}=1}
    \innerproduct{\mathbf{x}}{H\mathbf{x}}                              \\
    &=
    \lambda_1(A)+\lambda_1(H).
\end{align*}
Hence
\[
    \lambda_1(A+H)-\lambda_1(A)
    \le
    \lambda_1(H).
\]
\end{pfbox}


\begin{proposition}
For Hermitian matrices $A,H\in M_n(\mathbb{C})$,
\[
    \lambda_n(A+H)-\lambda_n(A)\le\lambda_1(H).
\]
\end{proposition}

\begin{pfbox}
For $\norm{\mathbf{x}}=1$, write
$f_A(\mathbf{x})=\innerproduct{\mathbf{x}}{A\mathbf{x}}$ and
$f_H(\mathbf{x})=\innerproduct{\mathbf{x}}{H\mathbf{x}}$.
By the Rayleigh quotient characterization,
\[
    \lambda_n(A+H)
    =
    \min_{\norm{\mathbf{x}}=1}
    \bp{f_A(\mathbf{x})+f_H(\mathbf{x})}.
\]
Since $f_H(\mathbf{x})\le \lambda_1(H)$ for every $\norm{\mathbf{x}}=1$, we have
\begin{align*}
    \lambda_n(A+H)
    &=
    \min_{\norm{\mathbf{x}}=1}
    \bp{f_A(\mathbf{x})+f_H(\mathbf{x})}                                \\
    &\le
    \min_{\norm{\mathbf{x}}=1}
    \bp{f_A(\mathbf{x})+\lambda_1(H)}                                    \\
    &=
    \min_{\norm{\mathbf{x}}=1} f_A(\mathbf{x})
    +
    \lambda_1(H)                                                         \\
    &=
    \lambda_n(A)+\lambda_1(H).
\end{align*}
Therefore,
\[
    \lambda_n(A+H)-\lambda_n(A)
    \le
    \lambda_1(H).
\]
\end{pfbox}


\begin{proposition}
For Hermitian matrices $A,H\in M_n(\mathbb{C})$,
\[
    \lambda_1(A+H)-\lambda_1(A)\ge\lambda_n(H).
\]
\end{proposition}

\begin{pfbox}
Apply the previous proposition to $-A$ and $-H$:
\[
    \lambda_n((-A)+(-H))-\lambda_n(-A)
    \le
    \lambda_1(-H).
\]
Since
\[
    \lambda_n(-(A+H))=-\lambda_1(A+H),
    \quad
    \lambda_n(-A)=-\lambda_1(A),
    \quad
    \lambda_1(-H)=-\lambda_n(H),
\]
we obtain
\[
    -\lambda_1(A+H)+\lambda_1(A)
    \le
    -\lambda_n(H).
\]
Thus
\[
    \lambda_1(A+H)-\lambda_1(A)
    \ge
    \lambda_n(H).
\]
Combining this with Proposition~\ref{prop:first-eigenvalue-diff-upper-bound}, we get
\[
    \lambda_n(H)
    \le
    \lambda_1(A+H)-\lambda_1(A)
    \le
    \lambda_1(H).
\]
\end{pfbox}


\begin{theorem}[Weyl Inequality: Simple Version]
Let $A,H\in M_n(\mathbb{C})$ be Hermitian. For each $i=1,2,\dots,n$,
\[
    \lambda_i(A+H)-\lambda_i(A)
    \le
    \lambda_1(H).
\]
\end{theorem}

\begin{pfbox}
By Courant--Fischer,
\[
    \lambda_i(A+H)
    =
    \max_{\substack{E\subseteq \mathbb{C}^n \\ \dim E=i}}
    \min_{\substack{\mathbf{x}\in E \\ \norm{\mathbf{x}}=1}}
    \innerproduct{\mathbf{x}}{(A+H)\mathbf{x}}.
\]
We want to show
\[
    \max_E
    \min_{\substack{\mathbf{x}\in E \\ \norm{\mathbf{x}}=1}}
    \innerproduct{\mathbf{x}}{(A+H)\mathbf{x}}
    \le
    \lambda_i(A)+\lambda_1(H).
\]
It is enough to show that for every subspace $E\subseteq\mathbb{C}^n$ with
$\dim E=i$,
\[
    \min_{\substack{\mathbf{x}\in E \\ \norm{\mathbf{x}}=1}}
    \innerproduct{\mathbf{x}}{(A+H)\mathbf{x}}
    \le
    \lambda_i(A)+\lambda_1(H).
\]

Choose a subspace $W_A\subseteq\mathbb{C}^n$ with $\dim W_A=n-i+1$ such that
\[
    \max_{\substack{\mathbf{y}\in W_A \\ \norm{\mathbf{y}}=1}}
    \innerproduct{\mathbf{y}}{A\mathbf{y}}
    =
    \lambda_i(A).
\]
For any subspace $E$ with $\dim E=i$, we have
$\dim E+\dim W_A=i+(n-i+1)=n+1$. Hence $E\cap W_A\neq\{\mathbf{0}\}$.

Choose a unit vector $\mathbf{x}\in E\cap W_A$. Then
\[
    \innerproduct{\mathbf{x}}{A\mathbf{x}}
    \le
    \lambda_i(A),
    \qquad
    \innerproduct{\mathbf{x}}{H\mathbf{x}}
    \le
    \lambda_1(H).
\]
Hence
\[
    \innerproduct{\mathbf{x}}{(A+H)\mathbf{x}}
    =
    \innerproduct{\mathbf{x}}{A\mathbf{x}}
    +
    \innerproduct{\mathbf{x}}{H\mathbf{x}}
    \le
    \lambda_i(A)+\lambda_1(H).
\]
Since this is true for some unit vector $\mathbf{x}\in E$,
\[
    \min_{\substack{\mathbf{z}\in E \\ \norm{\mathbf{z}}=1}}
    \innerproduct{\mathbf{z}}{(A+H)\mathbf{z}}
    \le
    \lambda_i(A)+\lambda_1(H).
\]
Taking the maximum over all such $E$ gives the desired inequality.
\end{pfbox}


\begin{corollary}
For Hermitian matrices $A,H\in M_n(\mathbb{C})$,
\[
    \lambda_i(A+H)
    \ge
    \lambda_i(A)+\lambda_n(H).
\]
Consequently,
\[
    \lambda_n(H)
    \le
    \lambda_i(A+H)-\lambda_i(A)
    \le
    \lambda_1(H).
\]
\end{corollary}

\begin{pfbox}
Use
\[
    \lambda_i(A+H)
    =
    -\lambda_{n-i+1}(-(A+H)).
\]
By the previous theorem applied to $-A$ and $-H$,
\[
    \lambda_{n-i+1}((-A)+(-H))
    \le
    \lambda_{n-i+1}(-A)+\lambda_1(-H).
\]
Multiplying by $-1$ gives
\begin{align*}
    \lambda_i(A+H)
    &\ge
    -\lambda_{n-i+1}(-A)-\lambda_1(-H)                                  \\
    &=
    \lambda_i(A)+\lambda_n(H).
\end{align*}
Combining this with the upper bound gives
\[
    \lambda_n(H)
    \le
    \lambda_i(A+H)-\lambda_i(A)
    \le
    \lambda_1(H).
\]
\end{pfbox}



\begin{theorem}[Continuity of eigenvalues]
For Hermitian matrices $A,B\in M_n(\mathbb{C})$,
\[
    \abs{\lambda_i(A)-\lambda_i(B)}
    \le
    \norm{A-B}_F,
    \qquad
    i=1,2,\dots,n.
\]
Hence the mapping from a Hermitian matrix to its $i$-th eigenvalue is
$1$-Lipschitz continuous with respect to the Frobenius norm, and hence continuous.
\end{theorem}

\begin{pfbox}
Apply the previous corollary with base matrix $B$ and perturbation
$H=A-B$. Then $B+H=A$, so
\[
    \lambda_n(A-B)
    \le
    \lambda_i(A)-\lambda_i(B)
    \le
    \lambda_1(A-B).
\]
Therefore,
\[
    \abs{\lambda_i(A)-\lambda_i(B)}
    \le
    \max\{\abs{\lambda_1(A-B)},\abs{\lambda_n(A-B)}\}.
\]

Using $\abs{\lambda_j(C)}\le\norm{C}_F$ for every $j$ and taking $C=A-B$, we get
\[
    \max\{\abs{\lambda_1(A-B)},\abs{\lambda_n(A-B)}\}
    \le
    \norm{A-B}_F.
\]
Therefore,
\[
    \abs{\lambda_i(A)-\lambda_i(B)}
    \le
    \norm{A-B}_F.
\]
\end{pfbox}




\begin{theorem}
For any $C\in M_n(\mathbb{C})$ and every eigenvalue $\lambda_j(C)$ of $C$, $\abs{\lambda_j(C)}\le\norm{C}_F$.
\end{theorem}


\begin{pfbox}
Let $\rho(C)=\max_j\abs{\lambda_j(C)}$ be the spectral radius of $C$.
It suffices to show that $\rho(C)\le\norm{C}_F$.

By JCF, there exists an invertible matrix $P$ such that $C=PJP^{-1}$,
where $J$ is a Jordan matrix. Hence $C^m=PJ^mP^{-1}$ for every $m\ge 1$.

We first show that
\[
    \lim_{m\to\infty}\norm{C^m}_F^{1/m}
    =
    \rho(C).
\]
It is enough to prove this for $J$. 
By submultiplicativity of the Frobenius norm,
\[
    \norm{C^m}_F
    =
    \norm{PJ^mP^{-1}}_F
    \le
    \norm{P}_F\norm{J^m}_F\norm{P^{-1}}_F.
\]
Let $K=\norm{P}_F\norm{P^{-1}}_F$. Then
$\norm{C^m}_F\le K\norm{J^m}_F$.

On the other hand,
\[
    \norm{J^m}_F
    =
    \norm{P^{-1}C^mP}_F
    \le
    \norm{P^{-1}}_F\norm{C^m}_F\norm{P}_F
    =
    K\norm{C^m}_F.
\]
Therefore,
\[
    \frac{1}{K}\norm{J^m}_F
    \le
    \norm{C^m}_F
    \le
    K\norm{J^m}_F.
\]
Taking the $m$-th root gives
\[
    K^{-1/m}\norm{J^m}_F^{1/m}
    \le
    \norm{C^m}_F^{1/m}
    \le
    K^{1/m}\norm{J^m}_F^{1/m}.
\]
Since $K>0$ is independent of $m$, we have
$K^{1/m}=e^{(\log K)/m}\to 1$ and
$K^{-1/m}=e^{-(\log K)/m}\to 1$. Thus $C^m$ and $J^m$ have the same
growth rate after taking the $m$-th root.

Now consider one Jordan block $J_r(\lambda)=\lambda I+N$, where $N^r=0$.

If $\lambda=0$, then $J_r(0)=N$ is nilpotent, so $J_r(0)^m=0$ for all
$m\ge r$. Hence
\[
    \lim_{m\to\infty}
    \norm{J_r(0)^m}_F^{1/m}
    =
    0
    =
    \abs{\lambda}.
\]

Now assume $\lambda\neq 0$. By the binomial theorem,
\[
    J_r(\lambda)^m
    =
    (\lambda I+N)^m
    =
    \sum_{q=0}^{r-1}
    \binom{m}{q}\lambda^{m-q}N^q.
\]
The sum stops at $r-1$ because $N^r=0$. Since $q\le r-1$ is bounded,
$\binom{m}{q}$ grows at most polynomially in $m$. Hence there exists a
constant $C_0>0$ such that
\[
    \norm{J_r(\lambda)^m}_F
    \le
    C_0m^{r-1}\abs{\lambda}^m
\]
for all sufficiently large $m$. Taking the $m$-th root gives
\[
    \norm{J_r(\lambda)^m}_F^{1/m}
    \le
    C_0^{1/m}(m^{r-1})^{1/m}\abs{\lambda}.
\]
Since $C_0^{1/m}\to 1$ and
$(m^{r-1})^{1/m}=e^{((r-1)\log m)/m}\to 1$, we obtain
\[
    \limsup_{m\to\infty}
    \norm{J_r(\lambda)^m}_F^{1/m}
    \le
    \abs{\lambda}.
\]

On the other hand, the diagonal entries of $J_r(\lambda)^m$ are all
$\lambda^m$. Since the Frobenius norm is at least the absolute value of any
entry, we have $\norm{J_r(\lambda)^m}_F\ge \abs{\lambda}^m$. Taking the
$m$-th root gives $\norm{J_r(\lambda)^m}_F^{1/m}\ge \abs{\lambda}$, and hence
\[
    \liminf_{m\to\infty}
    \norm{J_r(\lambda)^m}_F^{1/m}
    \ge
    \abs{\lambda}.
\]
Combining the upper and lower estimates,
\[
    \abs{\lambda}
    \le
    \liminf_{m\to\infty}
    \norm{J_r(\lambda)^m}_F^{1/m}
    \le
    \limsup_{m\to\infty}
    \norm{J_r(\lambda)^m}_F^{1/m}
    \le
    \abs{\lambda}.
\]
Therefore,
\[
    \lim_{m\to\infty}
    \norm{J_r(\lambda)^m}_F^{1/m}
    =
    \abs{\lambda}.
\]

Now write $J$ as a block diagonal matrix
\[
    J
    =
    \begin{pmatrix}
        J_{r_1}(\lambda_1) &        & \mathbf{O} \\
                           & \ddots &             \\
        \mathbf{O}         &        & J_{r_s}(\lambda_s)
    \end{pmatrix}.
\]
Then
\[
    \norm{J^m}_F^2
    =
    \sum_{\ell=1}^s
    \norm{J_{r_\ell}(\lambda_\ell)^m}_F^2.
\]
For each block, we have shown that
$\norm{J_{r_\ell}(\lambda_\ell)^m}_F^{1/m}\to \abs{\lambda_\ell}$.

Let $M=\max_{1\le \ell\le s}\abs{\lambda_\ell}$. Then $M=\rho(C)$.
For the lower bound, choose $\ell_0$ such that $\abs{\lambda_{\ell_0}}=M$.
Since $\norm{J^m}_F\ge \norm{J_{r_{\ell_0}}(\lambda_{\ell_0})^m}_F$, we get
\[
    \liminf_{m\to\infty}
    \norm{J^m}_F^{1/m}
    \ge
    M.
\]
For the upper bound, fix $\varepsilon>0$. For each block, we have
\[
    \lim_{m\to\infty}
    \norm{J_{r_\ell}(\lambda_\ell)^m}_F^{1/m}
    =
    \abs{\lambda_\ell}
    \le
    M.
\]
Therefore, for each $\ell$, there exists $N_\ell$ such that whenever
$m\ge N_\ell$,
\[
    \norm{J_{r_\ell}(\lambda_\ell)^m}_F^{1/m}
    \le
    M+\varepsilon.
\]
Equivalently,
\[
    \norm{J_{r_\ell}(\lambda_\ell)^m}_F
    \le
    (M+\varepsilon)^m.
\]
Since there are only finitely many blocks, we can choose
$N=\max\{N_1,\ldots,N_s\}$. Then for all $m\ge N$ and all
$\ell=1,\ldots,s$,
\[
    \norm{J_{r_\ell}(\lambda_\ell)^m}_F
    \le
    (M+\varepsilon)^m.
\]

Since $J$ is block diagonal,
\[
    \norm{J^m}_F^2
    =
    \sum_{\ell=1}^s
    \norm{J_{r_\ell}(\lambda_\ell)^m}_F^2.
\]
Thus, for all $m\ge N$,
\[
    \norm{J^m}_F^2
    \le
    \sum_{\ell=1}^s (M+\varepsilon)^{2m}
    =
    s(M+\varepsilon)^{2m}.
\]
Taking square roots gives
\[
    \norm{J^m}_F
    \le
    \sqrt{s}(M+\varepsilon)^m.
\]
Taking the $m$-th root gives
\[
    \norm{J^m}_F^{1/m}
    \le
    s^{1/(2m)}(M+\varepsilon).
\]
Since $s^{1/(2m)}\to 1$, we obtain
\[
    \limsup_{m\to\infty}
    \norm{J^m}_F^{1/m}
    \le
    M+\varepsilon.
\]
Since $\varepsilon>0$ is arbitrary,
\[
    \limsup_{m\to\infty}
    \norm{J^m}_F^{1/m}
    \le
    M.
\]
Therefore,
\[
    \lim_{m\to\infty}
    \norm{J^m}_F^{1/m}
    =
    M
    =
    \rho(C).
\]
By the comparison between $C^m$ and $J^m$ above, we also get
\[
    \lim_{m\to\infty}
    \norm{C^m}_F^{1/m}
    =
    \rho(C).
\]

Finally, since $\norm{C^m}_F\le \norm{C}_F^m$, we take the $m$-th root to give
$\norm{C^m}_F^{1/m}\le \norm{C}_F$. Letting $m\to\infty$, we obtain
$\rho(C)\le \norm{C}_F$.

Since $\abs{\lambda_j(C)}\le \rho(C)$ for every $j$, we conclude that
\[
    \abs{\lambda_j(C)}
    \le
    \norm{C}_F.
\]
\end{pfbox}

\begin{remark}
Now we show that if $C$ is Hermitian, then $\abs{\lambda_j(C)}\le\norm{C}_F$ for every $j$.

Since $C$ is Hermitian, by the Spectral Theorem, there exists a unitary matrix
$U$ such that $C=UDU^*$, where
$D=\operatorname{diag}\bigl(\lambda_1(C),\lambda_2(C),\ldots,\lambda_n(C)\bigr)$.
Hence
\begin{align*}
    \norm{C}_F^2
    &=\tr(C^*C) = \tr\bigl((UDU^*)^*(UDU^*)\bigr)\\
    &=\tr(UD^*DU^*)=\tr(U^*UD^*D)=\tr(D^*D)=\sum_{j=1}^n \abs{\lambda_j(C)}^2.
\end{align*}
Therefore, for every $j$,
\[
    \abs{\lambda_j(C)}^2
    \le
    \sum_{j=1}^n \abs{\lambda_\ell(C)}^2
    =
    \norm{C}_F^2.
\]
Thus
\[
    \abs{\lambda_j(C)}
    \le
    \norm{C}_F.
\]
\end{remark}





\begin{theorem}[Weyl's Inequality for Singular Values]
For general matrices $A,B\in M_{m\times n}(\mathbb{C})$,
\[
    \abs{\sigma_i(A)-\sigma_i(B)}
    \le
    \sqrt{2}\norm{A-B}_F.
\]
\end{theorem}

\begin{remark}
The mapping from a matrix to its singular values is therefore
$\sqrt{2}$-Lipschitz, and hence continuous.

A stronger version is also true:
\[
    \abs{\sigma_i(A)-\sigma_i(B)}
    \le
    \norm{A-B}_F,
\]
but this is harder to prove.
\end{remark}

\begin{pfbox}
We have already shown that for Hermitian matrices,
\[
    \abs{\lambda_i(A)-\lambda_i(B)}
    \le
    \norm{A-B}_F.
\]
For general matrices $A,B\in M_{m\times n}(\mathbb{C})$, consider the
Hermitian dilations
\[
    \widetilde{A}
    =
    \begin{pmatrix}
        \mathbf{O}_{m\times m} & A \\
        A^*                    & \mathbf{O}_{n\times n}
    \end{pmatrix},
    \qquad
    \widetilde{B}
    =
    \begin{pmatrix}
        \mathbf{O}_{m\times m} & B \\
        B^*                    & \mathbf{O}_{n\times n}
    \end{pmatrix}.
\]
Then $\widetilde{A}$ and $\widetilde{B}$ are Hermitian.

The goal is to compute the eigenvalues of $\widetilde{A}$. We have
\[
    \widetilde{A}^2
    =
    \begin{pmatrix}
        \mathbf{O} & A \\
        A^*        & \mathbf{O}
    \end{pmatrix}
    \begin{pmatrix}
        \mathbf{O} & A \\
        A^*        & \mathbf{O}
    \end{pmatrix}
    =
    \begin{pmatrix}
        AA^*      & \mathbf{O} \\
        \mathbf{O} & A^*A
    \end{pmatrix}.
\]
Hence the eigenvalues of $\widetilde{A}^2$ are the eigenvalues of $AA^*$
together with the eigenvalues of $A^*A$.

Let $\rank(A)=k$. Then the nonzero eigenvalues of $AA^*$ and $A^*A$ are
$\sigma_1(A)^2,\sigma_2(A)^2,\ldots,\sigma_k(A)^2$. Thus the nonzero
eigenvalues of $\widetilde{A}^2$ are
\[
    \sigma_1(A)^2,\ldots,\sigma_k(A)^2,
    \sigma_1(A)^2,\ldots,\sigma_k(A)^2.
\]
Therefore, every nonzero eigenvalue of $\widetilde{A}$ must be of the form
$\pm\sigma_i(A)$.

Now we show that the signs appear in pairs. Let
\[
    S
    =
    \begin{pmatrix}
        I_m        & \mathbf{O} \\
        \mathbf{O} & -I_n
    \end{pmatrix}.
\]
Then
\[
    S\widetilde{A}S
    =
    -\widetilde{A}.
\]
Equivalently, $\widetilde{A}S=-S\widetilde{A}$. Hence, if
$\widetilde{A}\mathbf{z}=\lambda\mathbf{z}$, then
\[
    \widetilde{A}(S\mathbf{z})
    =
    -S\widetilde{A}\mathbf{z}
    =
    -\lambda S\mathbf{z}.
\]
Since $S$ is invertible, $S\mathbf{z}\neq \mathbf{0}$. Therefore, whenever
$\lambda$ is an eigenvalue of $\widetilde{A}$, $-\lambda$ is also an
eigenvalue of $\widetilde{A}$ with the same multiplicity.

Therefore the eigenvalues of $\widetilde{A}$ are
\[
    \sigma_1(A),\ldots,\sigma_k(A),
    0,\ldots,0,
    -\sigma_k(A),\ldots,-\sigma_1(A).
\]
Here the zero eigenvalues fill the remaining dimensions. In particular, if
we pad the singular values by setting $\sigma_i(A)=0$ for
$i>\rank(A)$, then
\[
    \lambda_i(\widetilde{A})=\sigma_i(A),
    \qquad
    i=1,2,\ldots,\min\{m,n\}.
\]
Similarly,
\[
    \lambda_i(\widetilde{B})=\sigma_i(B),
    \qquad
    i=1,2,\ldots,\min\{m,n\}.
\]
Thus, by the continuity of eigenvalues for Hermitian matrices,
\[
    \abs{\sigma_i(A)-\sigma_i(B)}
    =
    \abs{\lambda_i(\widetilde{A})-\lambda_i(\widetilde{B})}
    \le
    \norm{\widetilde{A}-\widetilde{B}}_F.
\]

It remains to compute the Frobenius norm. Let $C=A-B$. Then
\[
    \widetilde{A}-\widetilde{B}
    =
    \begin{pmatrix}
        \mathbf{O} & C \\
        C^*        & \mathbf{O}
    \end{pmatrix}.
\]
Since $\widetilde{A}-\widetilde{B}$ is Hermitian,
\begin{align*}
    \norm{\widetilde{A}-\widetilde{B}}_F^2
    &=
    \tr\bp{(\widetilde{A}-\widetilde{B})^*
    (\widetilde{A}-\widetilde{B})} =  \tr\bp{(\widetilde{A}-\widetilde{B})^2}  \\                            
    &=
    \tr
    \begin{pmatrix}
        CC^*       & \mathbf{O} \\
        \mathbf{O} & C^*C
    \end{pmatrix}                                                        \\
    &=
    \tr(CC^*)+\tr(C^*C) = 2\tr(C^*C) = 2\norm{C}_F^2 = 2\norm{A-B}_F^2.
\end{align*}
Hence
\[
    \norm{\widetilde{A}-\widetilde{B}}_F
    =
    \sqrt{2}\norm{A-B}_F.
\]
Therefore,
\[
    \abs{\sigma_i(A)-\sigma_i(B)}
    \le
    \sqrt{2}\norm{A-B}_F,
    \qquad
    i=1,2,\ldots,\min\{m,n\}.
\]
\end{pfbox}
